%%%%%%%%%%%%%%%%%%%%%%%%%%%%%%%%%%%%%%%%%%%%%%%%%%%%%%%%%%%%%%%%%%%%%%%%%%%%%%%
%% Trabalho de Conclusão de Curso --- Engenharia de Software, iCEV
%%
%% Classe: abnTeX2 (ABNT NBR 14724, 6023, 6024, 6027, 6028, 6034, 10520),
%% no modelo de monografia do iCEV (monografia-tcc/main.tex). O modelo vem em
%% arquivo único; aqui ele está dividido, e este arquivo só monta o documento:
%%   preambulo/pacotes.tex      pacotes e caminhos para o repositório
%%   preambulo/estilo-icev.tex  layout do modelo: logo, capa, folha de rosto
%%   preambulo/dados.tex        título, autor, orientador, curso
%%   pre-textuais/              aprovação, resumos, siglas, símbolos
%%   capitulos/  apendices/     texto
%%
%% COMPILAR:   make          (ou: make clean && make)
%% O Makefile executa pdflatex -> bibtex -> makeglossaries -> makeindex ->
%% pdflatex x2, na ordem verificada, e grava o PDF em build/main.pdf.
%%
%% Este diretório vive dentro do repositório, em docs/TCC/, e consome dele,
%% sem copiar nada:
%%   figuras       ../../experiments/figures/output/   (regeneráveis)
%%   hardware      ../../hardware/
%% Por isso o projeto não compila sozinho no Overleaf.
%%%%%%%%%%%%%%%%%%%%%%%%%%%%%%%%%%%%%%%%%%%%%%%%%%%%%%%%%%%%%%%%%%%%%%%%%%%%%%%

\documentclass[
	12pt,
	openright,
	oneside,			% troque para twoside se o curso pedir frente e verso
	a4paper,
	chapter=TITLE,
	english,
	brazil
	]{abntex2}

%%%%%%%%%%%%%%%%%%%%%%%%%%%%%%%%%%%%%%%%%%%%%%%%%%%%%%%%%%%%%%%%%%%%%%%%%%%%%%%
%% Pacotes do documento
%%
%% Os do modelo iCEV (indentfirst, abntex2cite) e os que o conteúdo
%% exige (siunitx, glossaries, nomencl, booktabs). A ordem importa:
%%   glossaries depois do hyperref, que a abntex2 carrega.
%%%%%%%%%%%%%%%%%%%%%%%%%%%%%%%%%%%%%%%%%%%%%%%%%%%%%%%%%%%%%%%%%%%%%%%%%%%%%%%

%% --- Fontes e codificação --------------------------------------------------
\usepackage[T1]{fontenc}
\usepackage[utf8]{inputenc}
\usepackage{indentfirst}
%% Times no corpo, como no modelo. O modelo usa mathptmx; newtx compõe a mesma
%% Times com matemática completa (negrito, letras gregas), que os capítulos
%% usam. As figuras do pipeline (experiments/figures/scripts/theme_kaelix.R)
%% compõem em Times New Roman para casar com este corpo.
\usepackage{newtxtext,newtxmath}
\usepackage{microtype}

%% --- Pacotes de conteúdo ---------------------------------------------------
\usepackage{graphicx}
\usepackage{xcolor}      % marcas vermelhas \emandamento; substitui o color do modelo
\usepackage{booktabs}    % réguas de tabela
\usepackage{amsmath}
\usepackage{caption}
%% Título antes da ilustração ou da tabela (NBR 14724, 5.8 e 5.9).
\captionsetup{justification=justified,singlelinecheck=false}
\captionsetup[figure]{position=top}
\captionsetup[table]{position=top}
\usepackage{letltxmacro} % \LetLtxMacro para o \fonte de estilo-icev.tex
\usepackage{siunitx}
%% Separador decimal do português: vírgula. Sem isto, \num{0,5} sai "0.5".
%% Incerteza: \SI{75,0 \pm 3,5}{\percent} sai "(75,0 ± 3,5) %", com a unidade
%% valendo para os dois números; para "75,0 % ± 3,5 %", troque bracket por repeat.
\sisetup{
  output-decimal-marker      = {,},
  group-separator            = {\,},
  range-phrase               = {~a~},
  list-final-separator       = {~e~},
  per-mode                   = symbol,
  uncertainty-mode           = separate,
  separate-uncertainty-units = bracket
}

%% --- Citações (ABNT, sistema autor-data) -----------------------------------
\usepackage[alf,bibleftalign]{abntex2cite}
%% Estilo local, ajustado à NBR 6023:2025 e à NBR 10520:2023; as alterações
%% estão listadas no topo do arquivo. O abntex2-alf da distribuição segue as
%% normas de 2018 e 2002 e não imprime DOI.
\bibliographystyle{abntex2-alf-2025}
%% A NBR 6023 não usa < > em volta do endereço eletrônico desde 2018; o
%% abntex2cite ainda os acrescenta a todo \url.
\def\UrlLeft{}
\def\UrlRight{}

%% --- Listas de siglas e de símbolos ----------------------------------------
%% Sem as opções que põem as listas no sumário: elementos pré-textuais não podem
%% constar nele (NBR 6027, 6.3).
\usepackage[acronym,nomain,nonumberlist]{glossaries}
\makeglossaries
\usepackage{nomencl}
\makenomenclature

%% --- Caminhos para o repositório -------------------------------------------
%% Definidos uma única vez aqui; os capítulos têm fallback idêntico.
\newcommand{\kxfig}[1]{../../experiments/figures/output/#1}
\newcommand{\kxhw}[1]{../../#1}

%%%%%%%%%%%%%%%%%%%%%%%%%%%%%%%%%%%%%%%%%%%%%%%%%%%%%%%%%%%%%%%%%%%%%%%%%%%%%%%
%% Estilo do modelo de monografia do iCEV
%%
%% Transcrito de monografia-tcc/main.tex: fontes de título, logo no cabeçalho,
%% capa, folha de rosto, ambiente de resumo, citação nas referências e cores
%% dos links. Divergências deliberadas do modelo, todas pelas normas ABNT:
%%   - Referências sem indicativo numérico (o modelo as numerava como capítulo).
%%   - Texto preto, inclusive nos links (NBR 14724, 5.1); o modelo usa azul.
%%   - Títulos em corpo 12, graduados por caixa-alta e negrito, com o mesmo destaque
%%     no sumário e no texto (NBR 14724, 5.1 e 5.4; NBR 6027). Por isso as seções
%%     secundárias não usam caixa-alta (sem section=TITLE em main.tex).
%%   - Um espaço entre indicativo e título (NBR 14724, 5.2.2).
%%   - Margem inferior de 2 cm e número da página a 2 cm da borda superior (NBR 14724,
%%     5.1 e 5.3); o modelo chega a cerca de 1,5 cm e a 2,4 cm.
%%   - Introdução em caixa-alta no sumário, como os demais capítulos.
%%   - Sem título corrente no cabeçalho: o logo de 15 cm cobre quase toda a
%%     largura do texto (16 cm) e escondia o título atrás da imagem. A NBR 14724
%%     só exige a paginação no canto superior direito.
%%   - Introdução numerada, como seção primária (NBR 14724, 5.2.3; NBR 6024); o modelo a
%%     deixa sem número. Decisão do usuário, em 2026-11-01.
%%   - Referências alinhadas à esquerda, com título em negrito (NBR 6023).
%% Os dados (título, autor, curso) ficam em preambulo/dados.tex.
%%%%%%%%%%%%%%%%%%%%%%%%%%%%%%%%%%%%%%%%%%%%%%%%%%%%%%%%%%%%%%%%%%%%%%%%%%%%%%%

%% --- Títulos -----------------------------------------------------------------
\renewcommand{\ABNTEXchapterfont}{\rmfamily\bfseries}
\renewcommand{\ABNTEXchapterfontsize}{\normalsize}
\renewcommand{\ABNTEXsectionfontsize}{\normalsize}
\renewcommand{\ABNTEXsubsectionfont}{\rmfamily\mdseries}
\renewcommand{\ABNTEXsubsectionfontsize}{\normalsize}
\setsecnumformat{\csname the#1\endcsname\ }
%% Um espaço entre número e título (NBR 14724, 5.2.2): a definição da classe deixa dois.
\makeatletter
\renewcommand{\printchapternum}{%
	\tocprintchapter
	\setboolean{abntex@innonumchapter}{false}%
	\chapnumfont
	\ifthenelse{\boolean{abntex@apendiceousecao}}%
		{\space\thechapter\tocinnonumchapter\ABNTEXcaptiondelim}%
		{\thechapter\space}%
}
\makeatother
%% Título separado do texto por uma linha em branco de espaço 1,5 (NBR 14724, 5.2.2).
\setlength{\afterchapskip}{0pt}
%% Destaque do sumário igual ao do texto: terciária sem negrito (NBR 14724, 5.4).
\renewcommand{\cftsubsectionfont}{\normalfont\normalsize}
\renewcommand{\cftsubsectionpagefont}{\normalfont\normalsize}
%% Apêndice no sumário: "APÊNDICE A --- TÍTULO".
\cftinsertcode{A}{%
	\renewcommand{\cftchapteraftersnum}{~\textemdash~}%
	\renewcommand{\cftchapterpresnum}{\unskip}%% a classe deixa dois espaços após "APÊNDICE"
}
\renewcommand{\tocinnonumchapter}{%
	%% 2 em cobre "A ---"; o memoir soma a largura de "APÊNDICE", e a segunda linha alinha sob
	%% o início do título.
	\addtocontents{toc}{\cftsetindents{chapter}{\cftlastnumwidth}{2em}}%
	\cftinserthook{toc}{A}%
}
%% O sumário já recebe o espaço que antecede a primeira entrada.
\renewcommand{\aftertoctitle}{\par\vspace{6.3pt}}
\renewcommand{\afterloftitle}{\par\vspace{\baselineskip}}
\renewcommand{\afterlottitle}{\par\vspace{\baselineskip}}
\setbeforesecskip{\baselineskip}
\setaftersecskip{\baselineskip}
\setbeforesubsecskip{\baselineskip}
\setaftersubsecskip{\baselineskip}
%% Equações em numeração corrida, como figuras e tabelas.
\counterwithout{equation}{chapter}

%% --- Notas, ilustrações e tabelas ---------------------------------------------
%% Filete de 5 cm a partir da margem esquerda (NBR 14724, 5.2).
\renewcommand{\footnoterule}{\kern-3pt\hrule width 5cm\kern 2.6pt}
%% Tabela em corpo 12; só fonte, legenda e nota ficam menores (NBR 14724, 5.1).
\renewcommand{\IBGEtabfontsize}{\normalsize}
%% Com o título acima, o caption não deixa espaço antes da fonte.
\LetLtxMacro{\fonteabntex}{\fonte}
\renewcommand{\fonte}[2][\fontename]{\par\vspace{3pt}\fonteabntex[#1]{#2}}
%% Legenda abaixo da fonte (NBR 14724, 5.8); o título vem antes da ilustração.
\newcommand{\legenda}[1]{\nota[Legenda]{#1}}
%% Mesma autoria, datas diferentes: (Autor, 2009, 2022), como pede a NBR 10520, 6.1.
\newcommand{\citemesmoautor}[3]{(#1, \citeyear{#2}, \citeyear{#3})}

%% --- Referências -------------------------------------------------------------
%% Sem "Citado na página" do modelo: a NBR 6023 não prevê esse elemento.
\renewcommand{\bibsection}{%
	\chapter*{\bibname}%
	\addcontentsline{toc}{chapter}{\MakeTextUppercase{\bibname}}%
	\bibmark
	\vspace*{21.4pt}%% uma linha em branco de espaço 1,5 (NBR 14724, 5.2.2)
	\prebibhook
}

%% --- Cabeçalho com logo --------------------------------------------------------
\newcommand{\logoicevheader}{\includegraphics[width=15cm]{imagens/logo-icev.jpg}}

\makepagestyle{icevpretextual}
\makeevenhead{icevpretextual}{}{\logoicevheader}{}
\makeoddhead{icevpretextual}{}{\logoicevheader}{}
\aliaspagestyle{titlingpage}{icevpretextual}

\makepagestyle{icevheadings}
%% O número sobe para ficar a 2 cm da borda superior (NBR 14724, 5.3).
\newcommand{\icevnumeropagina}{\raisebox{0.35cm}{\ABNTEXfontereduzida\thepage}}
\makeevenhead{icevheadings}{\icevnumeropagina}{\logoicevheader}{}
\makeoddhead{icevheadings}{}{\logoicevheader}{\icevnumeropagina}
\makeheadrule{icevheadings}{\textwidth}{\normalrulethickness}

\makepagestyle{icevchapfirst}
\makeoddhead{icevchapfirst}{}{\logoicevheader}{\icevnumeropagina}

%% O modelo usa 2.1cm, mas o logo mede 60,6pt e não cabe (aviso do memoir e
%% \vbox overfull em toda página). Com 2.15cm cabe, e cabeçalho mais
%% \headsep continuam dentro da margem superior de 3 cm.
\setlength{\headheight}{2.15cm}
\setlength{\headsep}{0.5cm}
\checkandfixthelayout
%% Sem este ajuste, a última linha de uma página cheia fica a cerca de 1,7 cm da borda;
%% a margem inferior é de 2 cm (NBR 14724, 5.1).
\addtolength{\textheight}{-0.5cm}
\checkandfixthelayout[fixed]

%% --- Links ---------------------------------------------------------------------
%% pdftitle e pdfauthor ficam em dados.tex: aqui \@title e \@author ainda estão vazios.
\makeatletter
\hypersetup{
	pdfsubject  = {Trabalho de Conclusão de Curso},
	pdfcreator  = {LaTeX with abnTeX2},
	pdfkeywords = {manutenção preditiva; análise de vibração; detecção de anomalias;
	               sistemas embarcados; instrumentação de baixo custo},
	hidelinks,
	bookmarksdepth = 4
}
\makeatother

%% --- Parágrafo -----------------------------------------------------------------
\setlength{\parindent}{1.25cm}
\setlength{\parskip}{0pt}

%% --- Capa ----------------------------------------------------------------------
\renewcommand{\imprimircapa}{%
	\begin{capa}
		\centering
		{\ABNTEXchapterfont\bfseries\instituicaoicev\par}
		{\ABNTEXchapterfont\bfseries\cursoicev\par}

		\vfill
		{\ABNTEXchapterfont\bfseries\imprimirautor}
		\vfill

		\begin{center}
			\ABNTEXchapterfont\bfseries\imprimirtitulo
		\end{center}
		\vfill

		{\bfseries\imprimirlocal\par}
		{\bfseries\imprimirdata}
		\vspace*{1cm}
	\end{capa}
}

%% --- Natureza do trabalho --------------------------------------------------------
%% Do meio da mancha até a margem direita, em espaço simples, na folha de rosto e na
%% de aprovação (NBR 14724, 5.2).
\newcommand{\icevnatureza}{%
	{\raggedleft
	\begin{minipage}{.5\textwidth}
		\SingleSpacing
		%% Espaços sem encolher: nesta coluna estreita, a linha apertada colava
		%% as palavras ao copiar do PDF. Sem hifenizar o nome da instituição.
		\spaceskip=\fontdimen2\font plus \fontdimen2\font\relax
		\hyphenpenalty=10000
		\imprimirpreambulo
	\end{minipage}\par}%
}

%% --- Folha de rosto ------------------------------------------------------------
\makeatletter
\renewcommand{\folhaderostocontent}{
	\begin{center}

		{\ABNTEXchapterfont\bfseries\imprimirautor}

		\vspace*{\fill}\vspace*{\fill}
		\begin{center}
			\ABNTEXchapterfont\bfseries\imprimirtitulo
		\end{center}
		\vspace*{\fill}

		\abntex@ifnotempty{\imprimirpreambulo}{%
			\icevnatureza
			\vspace*{\fill}
		}%

		{\abntex@ifnotempty{\imprimirinstituicao}{\imprimirinstituicao\vspace*{\fill}}}

		{\imprimirorientadorRotulo~\imprimirorientador\par}
		\abntex@ifnotempty{\imprimircoorientador}{%
			{\imprimircoorientadorRotulo~\imprimircoorientador}%
		}%
		\vspace*{\fill}

		{\bfseries\imprimirlocal}
		\par
		{\bfseries\imprimirdata}
		\vspace*{1cm}

	\end{center}
}
\makeatother

%% --- Resumo --------------------------------------------------------------------
\renewenvironment{resumo}[1][\resumoname]{%
	\PRIVATEbookmarkthis{#1}
	\renewcommand{\abstractnamefont}{\ABNTEXchapterfont\bfseries}
	\renewcommand{\abstractname}{\ABNTEXchapterupperifneeded{#1}}
	\begin{abstract}
}{\end{abstract}\PRIVATEclearpageifneeded}
\setlength{\absparsep}{18pt}
\addtolength{\abstitleskip}{1.7pt}

%% --- Trechos que mudam com as medições do dispositivo --------------------------
%% Definido aqui porque o resumo e o abstract também usam, antes das seções.
\providecommand{\emandamento}[1]{\textcolor{red}{#1}}

%% --- Quadros --------------------------------------------------------------------
%% Quadro é ilustração (NBR 14724, 5.8), com lista própria antes da de tabelas.
\newcommand{\quadroname}{Quadro}
\newcommand{\listofquadrosname}{Lista de quadros}
\newfloat[chapter]{quadro}{loq}{\quadroname}
\newlistof{listofquadros}{loq}{\listofquadrosname}
\newlistentry{quadro}{loq}{0}
\counterwithout{quadro}{chapter}
\captionsetup[quadro]{position=top}
\renewcommand{\cftquadroname}{\quadroname\space}
\renewcommand*{\cftquadroaftersnum}{\hfill\textendash\hfill}
\renewcommand{\afterloqtitle}{\par\vspace{\baselineskip}}

%% --- Referências -------------------------------------------------------------
%% Separadas por uma linha em branco de espaço simples (NBR 6023, 6.3).
\setlength{\bibitemsep}{10pt}
%% Data provável, entre colchetes (NBR 6023, 8.6): os colchetes no campo year quebram o
%% argumento opcional de \bibitem, por isso ela fica em uma macro.
\newcommand{\anomafaulda}{[201-?]}

%% --- Listas de siglas e de símbolos ------------------------------------------
%% Mesma separação entre título e texto das demais listas.
\renewcommand*{\glossarypreamble}{\vspace*{\baselineskip}}
\renewcommand{\nompreamble}{\vspace*{20pt}}

%%%%%%%%%%%%%%%%%%%%%%%%%%%%%%%%%%%%%%%%%%%%%%%%%%%%%%%%%%%%%%%%%%%%%%%%%%%%%%%
%% Dados do trabalho
%%%%%%%%%%%%%%%%%%%%%%%%%%%%%%%%%%%%%%%%%%%%%%%%%%%%%%%%%%%%%%%%%%%%%%%%%%%%%%%

%% TÍTULO PROVISÓRIO (working title). Finalizar só após o ensaio do protótipo:
%% o título deve refletir o que foi realizado, e os resultados ainda podem mudar.
%%   - Descritivo, sem afirmar resultado ("limites" entra só depois do ensaio).
%%   - Sem nome de projeto (Kaelix) nem modelo comercial (MPU6050, ESP32-S3):
%%     a NBR 14724 exige que o título permita indexação e recuperação.
%%   - Sem sigla (MEMS): desaconselhada em título por IEEE e Nature.
%%   - O subtítulo descreve o conteúdo; na versão final, é onde entra a
%%     contribuição apurada.
%% A quebra depois dos dois-pontos separa título e subtítulo nas três folhas; nos
%% metadados do PDF ela vira espaço.
\titulo{Monitoramento de vibração de baixo custo para manutenção preditiva:\texorpdfstring{\\}{ }aquisição e detecção embarcada de anomalias em motores}
\autor{Lauan Matheus da Rocha Alves}
\local{Teresina -- PI}
\data{2026}
\orientador{Lia Mariana Nascimento Brito}
%\coorientador{Nome do coorientador}
\preambulo{Trabalho de Conclusão de Curso apresentado como requisito parcial
para obtenção do título de Bacharel em Engenharia de Software do Instituto de
Ensino Superior -- iCEV.}

%% Metadados do PDF: precisam vir depois de \titulo e \autor.
\makeatletter
\hypersetup{pdftitle = {\@title}, pdfauthor = {\@author}}
\makeatother

%% Usados na capa do modelo iCEV (preambulo/estilo-icev.tex).
\newcommand{\cursoicev}{Curso de Bacharelado em Engenharia de Software}
\newcommand{\instituicaoicev}{Instituto de Ensino Superior -- iCEV}


\makeindex

%% Listas de siglas e símbolos (definidas antes do documento)
%%%%%%%%%%%%%%%%%%%%%%%%%%%%%%%%%%%%%%%%%%%%%%%%%%%%%%%%%%%%%%%%%%%%%%%%%%%%%%%
%% Lista de Siglas e Abreviaturas (pacote glossaries)
%%
%% Contém apenas as siglas efetivamente usadas nas seções 1 a 5 e no apêndice.
%% Acrescentar entrada aqui ao introduzir sigla nova no texto.
%%%%%%%%%%%%%%%%%%%%%%%%%%%%%%%%%%%%%%%%%%%%%%%%%%%%%%%%%%%%%%%%%%%%%%%%%%%%%%%

%% Instrumentação
\newacronym{MEMS}{MEMS}{sistemas microeletromecânicos (\emph{Micro-Electro-Mechanical Systems})}
\newacronym{IMU}{IMU}{unidade de medição inercial (\emph{Inertial Measurement Unit})}
\newacronym{BPFO}{BPFO}{frequência de passagem de esferas na pista externa (\emph{Ball Pass Frequency, Outer race})}
\newacronym{BPFI}{BPFI}{frequência de passagem de esferas na pista interna (\emph{Ball Pass Frequency, Inner race})}
\newacronym{DLPF}{DLPF}{filtro passa-baixas digital (\emph{Digital Low-Pass Filter})}

%% Normalização
\newacronym{ISO}{ISO}{Organização Internacional de Normalização (\emph{International Organization for Standardization})}

%% Detecção e avaliação
\newacronym{ROC}{ROC}{característica de operação do receptor (\emph{Receiver Operating Characteristic})}
\newacronym{AUC}{AUC}{área sob a curva (\emph{Area Under the Curve})}
\newacronym{pAUC}{pAUC}{área parcial sob a curva, restrita a uma faixa de falso alarme}
\newacronym{PRAUC}{PR-AUC}{área sob a curva de precisão e revocação}
\newacronym{TPR}{TPR}{taxa de detecção (\emph{True Positive Rate})}
\newacronym{DFT}{DFT}{transformada discreta de Fourier (\emph{Discrete Fourier Transform})}
\newacronym{FFT}{FFT}{transformada rápida de Fourier (\emph{Fast Fourier Transform})}
\newacronym{FPR}{FPR}{taxa de falso alarme (\emph{False Positive Rate})}

%% Computação embarcada e comunicação
\newacronym{TinyML}{TinyML}{aprendizado de máquina em dispositivos de recursos limitados}
\newacronym{LoRa}{LoRa}{modulação de longo alcance e baixa potência (\emph{Long Range})}
\newacronym{SF}{SF}{fator de espalhamento (\emph{Spreading Factor})}
\newacronym{CRC}{CRC}{verificação cíclica de redundância (\emph{Cyclic Redundancy Check})}
\newacronym{SRAM}{SRAM}{memória estática de acesso aleatório (\emph{Static Random-Access Memory})}
\newacronym{I2C}{I2C}{circuito interintegrado (\emph{Inter-Integrated Circuit})}
\newacronym{SPI}{SPI}{interface serial periférica (\emph{Serial Peripheral Interface})}
\newacronym{USB}{USB}{barramento serial universal (\emph{Universal Serial Bus})}
\newacronym{FIFO}{FIFO}{memória em que o primeiro dado a entrar é o primeiro a sair (\emph{First In, First Out})}
\newacronym{NaN}{NaN}{valor que indica ausência de número (\emph{Not a Number})}

%% Taxas de decisão, descritores e parâmetros do enlace, usados em forma literal no texto.
\newacronym{BW}{BW}{largura de banda (\emph{Bandwidth})}
\newacronym{CR}{CR}{taxa de codificação (\emph{Coding Rate})}
\newacronym{FC}{FC}{fator de crista}
\newacronym{FN}{FN}{falso negativo}
\newacronym{FP}{FP}{falso positivo}
\newacronym{VN}{VN}{verdadeiro negativo}
\newacronym{VP}{VP}{verdadeiro positivo}

%%%%%%%%%%%%%%%%%%%%%%%%%%%%%%%%%%%%%%%%%%%%%%%%%%%%%%%%%%%%%%%%%%%%%%%%%%%%%%%
%% O pacote glossaries só imprime as siglas EFETIVAMENTE USADAS via \gls{...}.
%% Como o texto dos capítulos usa as siglas em forma literal (e não \gls{}),
%% \glsaddall força todas as entradas acima para a lista impressa.
%%
%% Se no futuro o texto passar a usar \gls{}, remova a linha abaixo para que a
%% lista volte a conter apenas o que é citado.
%%%%%%%%%%%%%%%%%%%%%%%%%%%%%%%%%%%%%%%%%%%%%%%%%%%%%%%%%%%%%%%%%%%%%%%%%%%%%%%
\glsaddall

%%%%%%%%%%%%%%%%%%%%%%%%%%%%%%%%%%%%%%%%%%%%%%%%%%%%%%%%%%%%%%%%%%%%%%%%%%%%%%%
%% Lista de Símbolos (pacote nomencl)
%%
%% Contém os símbolos que o texto define e usa. O prefixo numérico mantém a ordem de
%% primeira ocorrência, exigida pela NBR 14724; ao acrescentar um símbolo, escolha o
%% prefixo pela posição dele no texto. A macro \nomunit acrescenta a unidade alinhada à
%% direita.
%%%%%%%%%%%%%%%%%%%%%%%%%%%%%%%%%%%%%%%%%%%%%%%%%%%%%%%%%%%%%%%%%%%%%%%%%%%%%%%

\renewcommand{\nomname}{Lista de símbolos}

%----------------------------------------------
\newcommand{\nomunit}[1]{%
\renewcommand{\nomentryend}{\hspace*{\fill}#1}}
%----------------------------------------------

%% --- Vibração e descritores --------------------------------------------------
\nomenclature[001]{\(x(t)\)}{Deslocamento\nomunit{\si{\meter}}}
\nomenclature[002]{\(t\)}{Tempo\nomunit{\si{\second}}}
\nomenclature[003]{\(v(t)\)}{Velocidade\nomunit{\si{\meter\per\second}}}
\nomenclature[004]{\(a(t)\)}{Aceleração\nomunit{\si{\meter\per\second\squared}}}
\nomenclature[005]{\(f\)}{Frequência\nomunit{\si{\hertz}}}
\nomenclature[006]{\(X\)}{Amplitude de deslocamento\nomunit{\si{\meter}}}
\nomenclature[007]{\(\omega\)}{Frequência angular\nomunit{\si{\radian\per\second}}}
\nomenclature[008]{\(N\)}{Número de amostras na janela}
\nomenclature[009]{\(x_n\)}{Amostra de índice \(n\) da janela}
\nomenclature[010]{\(\bar{x}\)}{Média do sinal na janela}
\nomenclature[011]{\(\sigma\)}{Desvio-padrão do sinal na janela\nomunit{\si{\meter\per\second\squared}}}
\nomenclature[012]{\(x_{\mathrm{rms}}\)}{Valor eficaz do sinal na janela\nomunit{\si{\meter\per\second\squared}}}
\nomenclature[013]{\(\kappa\)}{Curtose do sinal na janela}
\nomenclature[014]{\(x_{\mathrm{pico}}\)}{Valor de pico do sinal na janela}

%% --- Rolamento e aquisição ---------------------------------------------------
\nomenclature[022]{\(f_r\)}{Frequência de rotação do eixo\nomunit{\si{\hertz}}}
\nomenclature[023]{\(N_b\)}{Número de esferas do rolamento}
\nomenclature[024]{\(d\)}{Diâmetro da esfera\nomunit{\si{\milli\meter}}}
\nomenclature[025]{\(D_p\)}{Diâmetro primitivo do rolamento\nomunit{\si{\milli\meter}}}
\nomenclature[026]{\(\theta\)}{Ângulo de contato do rolamento\nomunit{\si{\radian}}}
\nomenclature[027]{\(f_s\)}{Taxa de amostragem\nomunit{\si{\hertz}}}
\nomenclature[028]{\(f_{\mathrm{ap}}\)}{Frequência aparente sob subamostragem sem filtro\nomunit{\si{\hertz}}}
\nomenclature[029]{\(X_k\)}{Coeficiente de ordem \(k\) da DFT}
\nomenclature[030]{\(\Delta f\)}{Resolução em frequência da DFT\nomunit{\si{\hertz}}}
\nomenclature[031]{\(B\)}{Largura da banda de ruído\nomunit{\si{\hertz}}}

%% --- Detector e avaliação ----------------------------------------------------
\nomenclature[032]{\(h(x)\)}{Comprimento do caminho de isolamento de \(x\)}
\nomenclature[033]{\(\psi\)}{Tamanho da subamostra de cada árvore}
\nomenclature[034]{\(c(\psi)\)}{Comprimento médio de caminho para subamostra de \(\psi\) pontos}
\nomenclature[035]{\(s(x,\psi)\)}{Escore de anomalia do Isolation Forest}
\nomenclature[036]{\(\alpha\)}{Fração da operação normal acima do limiar}
\nomenclature[037]{\(F_1\)}{Medida que combina precisão e TPR}
\nomenclature[038]{\(\beta\)}{Limite superior de falso alarme da área parcial}
\nomenclature[0385]{\(A_\beta\)}{Área sob a curva ROC entre 0 e \(\beta\)}
\nomenclature[039]{\(\mathrm{pAUC}(\beta)\)}{Área parcial padronizada sob a curva ROC}

%% --- Enlace LoRa -------------------------------------------------------------
\nomenclature[040]{\(R_s\)}{Taxa de símbolos LoRa\nomunit{\si{\per\second}}}
\nomenclature[041]{\(T_s\)}{Duração de um símbolo LoRa\nomunit{\si{\second}}}
\nomenclature[042]{\(T_{\mathrm{pacote}}\)}{Duração de um pacote LoRa\nomunit{\si{\second}}}
\nomenclature[043]{\(n_{\mathrm{pre}}\)}{Número de símbolos do preâmbulo}

%% --- Modelo de sinal (seção de métodos) --------------------------------------
\nomenclature[047]{\(A_k\)}{Amplitude da \(k\)-ésima harmônica da rotação\nomunit{\si{\meter\per\second\squared}}}
\nomenclature[048]{\(\varphi_k\)}{Fase da \(k\)-ésima harmônica\nomunit{\si{\radian}}}
\nomenclature[049]{\(D\)}{Amplitude do impacto do defeito\nomunit{\si{\meter\per\second\squared}}}
\nomenclature[050]{\(\zeta_d\)}{Taxa de amortecimento do impacto\nomunit{\si{\per\second}}}
\nomenclature[051]{\(f_{\mathrm{res}}\)}{Frequência da ressonância estrutural excitada\nomunit{\si{\hertz}}}
\nomenclature[052]{\(t_i\)}{Instante do \(i\)-ésimo impacto\nomunit{\si{\second}}}
\nomenclature[053]{\(\varepsilon(t)\)}{Ruído gaussiano branco\nomunit{\si{\meter\per\second\squared}}}
\nomenclature[054]{\(u(\cdot)\)}{Degrau unitário}
\nomenclature[055]{\(\delta_i\)}{Desvio aleatório que modela o escorregamento das esferas\nomunit{\si{\second}}}


\begin{document}

\selectlanguage{brazil}
\frenchspacing
\pagestyle{icevpretextual}
\aliaspagestyle{chapter}{icevpretextual}

%%%%%%%%%%%%%%%%%%%%%%%%%%%%%%%%%%%%%%%%%%%%%%%%%%%%%%%%%%%%%%%%%%%%%%%%%%%%%%%
%% ELEMENTOS PRÉ-TEXTUAIS
%%%%%%%%%%%%%%%%%%%%%%%%%%%%%%%%%%%%%%%%%%%%%%%%%%%%%%%%%%%%%%%%%%%%%%%%%%%%%%%

\imprimircapa
\imprimirfolhaderosto

%% Ficha catalográfica: obrigatória na NBR 14724, 4.2.1.1.2; em formato eletrônico, vem
%% logo após a página de rosto. É elaborada pela biblioteca do iCEV: inserir aqui o PDF
%% recebido, na versão final.
%\begin{fichacatalografica}
%  \includepdf{fig_ficha_catalografica.pdf}
%\end{fichacatalografica}

%%%%%%%%%%%%%%%%%%%%%%%%%%%%%%%%%%%%%%%%%%%%%%%%%%%%%%%%%%%%%%%%%%%%%%%%%%%%%%%
%% Folha de aprovação --- modelo iCEV, com os membros da banca (NBR 14724).
%%
%% Depois da defesa, troque o ambiente abaixo pelo PDF da folha assinada:
%%   \usepackage{pdfpages}   (em preambulo/pacotes.tex)
%%   \begin{folhadeaprovacao}
%%     \includepdf{folha-aprovacao-assinada.pdf}
%%   \end{folhadeaprovacao}
%%%%%%%%%%%%%%%%%%%%%%%%%%%%%%%%%%%%%%%%%%%%%%%%%%%%%%%%%%%%%%%%%%%%%%%%%%%%%%%

\begin{folhadeaprovacao}
	\begin{center}
		{\ABNTEXchapterfont\bfseries\imprimirautor}
		\vspace*{\fill}\vspace*{\fill}
		\begin{center}
			\ABNTEXchapterfont\bfseries\imprimirtitulo
		\end{center}
		\vspace*{\fill}
		\icevnatureza
		\vspace*{\fill}
	\end{center}

	\noindent{\bfseries\imprimirlocal, \underline{\hspace{0.8cm}} de \underline{\hspace{2.5cm}} de 20\underline{\hspace{0.5cm}}}

	\assinatura{\bfseries\imprimirorientador \\ Professor Orientador \\ Titulação \\ Instituição}
	\assinatura{\bfseries Membro da banca \\ Titulação \\ Instituição}
	\assinatura{\bfseries Membro da banca \\ Titulação \\ Instituição}
\end{folhadeaprovacao}


%% Dedicatória, agradecimentos e epígrafe são opcionais (NBR 14724).
%\begin{agradecimentos}
%  ...
%\end{agradecimentos}

%%%%%%%%%%%%%%%%%%%%%%%%%%%%%%%%%%%%%%%%%%%%%%%%%%%%%%%%%%%%%%%%%%%%%%%%%%%%%%%
%% Resumo --- ABNT NBR 6028:2021
%%   parágrafo único; voz ativa; sem citações;
%%   palavras-chave separadas por ponto e vírgula, iniciais minúsculas,
%%   finalizadas por ponto.
%%
%% Ao alterar a resposta ao problema de pesquisa, alterar também as Conclusões e o
%% abstract.
%%%%%%%%%%%%%%%%%%%%%%%%%%%%%%%%%%%%%%%%%%%%%%%%%%%%%%%%%%%%%%%%%%%%%%%%%%%%%%%

\begin{resumo}
A análise de vibração permite decidir a manutenção de motores pelo estado medido da máquina,
mas o custo da instrumentação restringe seu uso às máquinas críticas e deixa a maior parte
dos motores de uma planta sem monitoramento. Nós sensores de baixo custo podem ampliar essa
cobertura, porém os trabalhos que os propõem raramente informam a banda e o piso de ruído do
sensor. Este trabalho investigou quais modos de falha uma cadeia de medição de baixo custo, com
acelerômetro MEMS amostrado a \SI{1}{\kilo\hertz}, microcontrolador ESP32 e enlace LoRa, permite
detectar, e quanto as decisões de estimação alteram o desempenho relatado. O objetivo foi medir a
detecção por modo de falha e calibrar o limiar de um detector Isolation Forest para uma taxa de
falso alarme fixada. Os experimentos usaram \num{1951} gravações de bancada do conjunto público
MAFAULDA, decimadas para \SI{1}{\kilo\hertz}, e sinais sintetizados, com a integração testada
contra um valor analítico. Em dado real, a \SI{2,0 \pm 1,3}{\percent} de falso alarme, a detecção
foi de \SI{75,0 \pm 3,5}{\percent} em desbalanceamento, de \SI{14,8}{\percent} e
\SI{20,7}{\percent} em defeito de rolamento sem massa sobreposta e de \SI{5,0}{\percent} e
\SI{7,2}{\percent} em desalinhamento; com o limiar padrão, o falso alarme foi de \SI{31,2 \pm
4,2}{\percent}. Em sinal sintetizado, a soma acumulada sem remoção da média errou a velocidade
eficaz em \SI{73}{\percent}, e em \SI{716}{\percent} com viés contido na especificação do sensor,
enquanto a integração no domínio da frequência ficou a menos de \SI{0,01}{\percent} do valor
analítico; a área parcial sob a curva ROC caiu até \SI{20,3}{\percent} entre as cadeias de
\SI{20}{\kilo\hertz} e \SI{1}{\kilo\hertz}, contra \SI{3,7}{\percent} da área global; e um
classificador supervisionado atingiu \SI{99,0}{\percent} de acurácia sob partição aleatória de
janelas e \SI{63,2}{\percent} sob partição por ensaio. \emandamento{Em duas montagens de bancada,
sem motor e sem falha induzida, o dispositivo amostrou a cerca de \SI{999,8}{\hertz} em três ciclos,
e o receptor aceitou \num{270} pacotes.} Os resultados indicam que, com esse detector e quatro
descritores, a cadeia permite detectar desbalanceamento sob taxa de falso alarme fixada e detecta
pouco defeito de rolamento sem desbalanceamento sobreposto. Indicam também que o desempenho de
dispositivos dessa classe só pode ser comparado quando a banda, o limiar e o protocolo de validação
são informados.

\vspace{\onelineskip}
\noindent
\textbf{Palavras-chave}: manutenção preditiva; análise de vibração; detecção de anomalias;
sistemas embarcados; instrumentação de baixo custo.
\end{resumo}

%%%%%%%%%%%%%%%%%%%%%%%%%%%%%%%%%%%%%%%%%%%%%%%%%%%%%%%%%%%%%%%%%%%%%%%%%%%%%%%
%% Abstract --- tradução do resumo.
%%%%%%%%%%%%%%%%%%%%%%%%%%%%%%%%%%%%%%%%%%%%%%%%%%%%%%%%%%%%%%%%%%%%%%%%%%%%%%%

\begin{resumo}[Abstract]
\begin{otherlanguage*}{english}
%% Ponto decimal só neste ambiente; o resto do documento usa vírgula.
\sisetup{output-decimal-marker = {.}}
Vibration analysis lets maintenance decisions follow the measured condition of a motor, but the
cost of instrumentation limits its use to critical machines and leaves most motors in a plant
unmonitored. Low-cost sensor nodes could extend this coverage, yet the studies proposing them rarely
report the bandwidth and noise floor of the sensor. This work investigated which motor failure modes
a low-cost measurement chain, made up of a MEMS accelerometer sampled at \SI{1}{\kilo\hertz}, an ESP32
microcontroller and a LoRa link, can detect, and how much estimation decisions change the reported
performance. The aim was to measure detection by failure mode and to calibrate the threshold of an
Isolation Forest detector for a fixed false alarm rate. The experiments used \num{1951} test-rig
recordings from the public MAFAULDA dataset, decimated to \SI{1}{\kilo\hertz}, and synthesised signals,
with integration tested against an analytical value. On real data, at a \SI{2.0 \pm 1.3}{\percent}
false alarm rate, detection was \SI{75.0 \pm 3.5}{\percent} for imbalance, \SI{14.8}{\percent} and
\SI{20.7}{\percent} for bearing faults with no added mass and \SI{5.0}{\percent} and \SI{7.2}{\percent}
for misalignment; with the default threshold, the false alarm rate was \SI{31.2 \pm 4.2}{\percent}. On
synthesised signals, a cumulative sum without mean removal missed the RMS velocity by \SI{73}{\percent},
and by \SI{716}{\percent} with a bias within the sensor specification, whereas frequency-domain
integration came within \SI{0.01}{\percent} of the analytical value; the partial area under the ROC curve
fell by up to \SI{20.3}{\percent} between the \SI{20}{\kilo\hertz} and \SI{1}{\kilo\hertz} chains, against
\SI{3.7}{\percent} for the global area; and a supervised classifier reached \SI{99.0}{\percent} accuracy
with random window partitioning and \SI{63.2}{\percent} with partitioning by test run. \emandamento{In two
bench setups, with no motor and no induced fault, the device sampled at about \SI{999.8}{\hertz} over three
cycles and the receiver accepted \num{270} packets.} The results indicate that, with this detector and four
features, the chain can detect imbalance at a fixed false alarm rate and detects few bearing faults when no
imbalance is added. They also indicate that the performance of devices of this class can only be compared
when the bandwidth, threshold and validation protocol are reported.

\vspace{\onelineskip}
\noindent
\textbf{Keywords}: predictive maintenance; vibration analysis; anomaly detection;
embedded systems; low-cost instrumentation.
\end{otherlanguage*}
\end{resumo}


%% Listas
\renewcommand{\listfigurename}{Lista de figuras}
\pdfbookmark[0]{\listfigurename}{lof}
\listoffigures*
\cleardoublepage

\pdfbookmark[0]{\listtablename}{lot}
\listoftables*
\cleardoublepage

\printglossary[type=\acronymtype, title={Lista de abreviaturas e siglas}, toctitle={Lista de abreviaturas e siglas}]
\cleardoublepage

\printnomenclature
\cleardoublepage

\pdfbookmark[0]{\contentsname}{toc}
\tableofcontents*
\cleardoublepage

%%%%%%%%%%%%%%%%%%%%%%%%%%%%%%%%%%%%%%%%%%%%%%%%%%%%%%%%%%%%%%%%%%%%%%%%%%%%%%%
%% ELEMENTOS TEXTUAIS
%%%%%%%%%%%%%%%%%%%%%%%%%%%%%%%%%%%%%%%%%%%%%%%%%%%%%%%%%%%%%%%%%%%%%%%%%%%%%%%
\textual
\pagestyle{icevheadings}
\aliaspagestyle{chapter}{icevchapfirst}

%%%%%%%%%%%%%%%%%%%%%%%%%%%%%%%%%%%%%%%%%%%%%%%%%%%%%%%%%%%%%%%%%%%%%%%%%%%%%%%
%% Capítulo 1 --- Introdução
%%
%% Citações: abntex2cite (sistema autor-data). \cite{} → (AUTOR, ano);
%%           \citeonline{} → Autor (ano).
%%
%% MARCAÇÃO DE TRECHO EM ANDAMENTO --------------------------------------------
%% \emandamento{...} destaca em vermelho os trechos que mudam com as medições do
%% dispositivo montado. É marcação de rascunho.
%%
%% Para REMOVER todas as marcas na versão final, troque a definição abaixo por:
%%     \newcommand{\emandamento}[1]{#1}
%% ou apague as chamadas \emandamento{...} do texto.
%%%%%%%%%%%%%%%%%%%%%%%%%%%%%%%%%%%%%%%%%%%%%%%%%%%%%%%%%%%%%%%%%%%%%%%%%%%%%%%

\makeatletter
\@ifundefined{emandamento}{%
  \@ifundefined{textcolor}%
    {\newcommand{\emandamento}[1]{#1}}%
    {\newcommand{\emandamento}[1]{\textcolor{red}{#1}}}%
}{}
\makeatother

%% Seção primária numerada: a NBR 14724 (5.2.3) e a NBR 6024 só admitem sem número os
%% títulos que listam, e a Introdução não está entre eles. O modelo do iCEV a deixa sem
%% número; a decisão de seguir a norma é do usuário (2026-11-01).
\chapter{Introdução}
\label{cap:introducao}

A vibração é a grandeza mais usada para acompanhar a condição de máquinas rotativas
\cite{randall2011livro}. A norma ISO 20816-3:2022, da International Organization for
Standardization (ISO), define como avaliá-la em máquinas industriais acima de
\SI{15}{\kilo\watt}, com rotação entre \SI{120}{r/min} e \SI{30000}{r/min}, e pede um
instrumento com resposta plana em \emph{ao menos}
\SIrange{10}{1000}{\hertz}~\cite{iso208163_2022}. Acelerômetros piezoelétricos atendem a
esse requisito, mas eles, seus condicionadores e os analisadores associados custam caro.
Por isso, o monitoramento costuma ficar restrito às máquinas consideradas críticas, e a
maior parte dos motores de uma planta continua sem instrumentação.

Uma alternativa estudada na literatura recente é o nó sensor de baixo custo: uma unidade
de medição inercial (IMU) baseada em sistemas microeletromecânicos (MEMS), ligada a um
microcontrolador que executa a detecção de anomalias e envia o resultado por rádio. O
algoritmo Isolation Forest~\cite{liu2008} já foi treinado e executado dentro de um
ESP32~\cite{antonini2023}, há variantes criadas para reduzir seu custo em aprendizado de
máquina para dispositivos de recursos limitados (TinyML)~\cite{barbariol2022}, e a
combinação de microcontrolador, acelerômetro MEMS de consumo e detecção não supervisionada
já foi aplicada à manutenção preditiva~\cite{kolok2025}. Esses trabalhos executam o detector
no microcontrolador. Eles não indicam quais falhas um sensor desse tipo consegue observar.

Duas etapas determinam o que essa medição entrega. A aquisição define que parte da
vibração chega ao processamento e depende da banda e do piso de ruído do sensor. A
estimação define o que se faz com as amostras: como integrar a aceleração, onde colocar o
limiar do detector, que métrica informar e como separar treino e teste. Uma componente que
não foi adquirida não pode ser recuperada depois.

Sobre a aquisição, os estudos revisados concordam entre si. Nenhum dos estudos de adequação de
sensor revisados testou dispositivo com limite superior abaixo de \SI{1500}{\hertz}~\cite{albarbar2008},
e os grupos que diagnosticam defeito de rolamento com MEMS escolhem sensores de dezenas de
quilohertz~\cite{jakobsen2024,landi2023}, porque a assinatura do defeito aparece modulada
em uma portadora de alta frequência~\cite{antoni2007,randall2011}. Sobre a estimação, há
menos trabalhos. O escore do Isolation Forest usa \num{0,5} como referência de
interpretação, sem calibração para os dados~\cite{liu2008}, e a implementação de
referência transforma esse valor em um deslocamento fixo, que não altera a ordem dos
escores~\cite{scikitlearn2026}. Os trabalhos que relatam bom desempenho com o detector
trocam o valor padrão, mas não testam a alternativa por varredura~\cite{antonini2023,brito2022}.
Entre os trabalhos que executam o detector em microcontrolador, a busca não encontrou
nenhum que variasse o parâmetro de contaminação sobre a mesma base e informasse a fração
da operação normal marcada como anomalia. Também não encontrou taxa de falso alarme
publicada para detector treinado no próprio dispositivo, que é o caso de uma máquina sem
histórico rotulado.

A validação é outro ponto fraco. Janelas consecutivas de um mesmo ensaio têm o mesmo nível
de ruído, o mesmo ganho do sensor e a mesma fixação. Quando essas janelas são sorteadas
entre treino e teste, o modelo pode reconhecer o ensaio em vez da falha.
\citeonline{varejao2025} chamam o fenômeno de \emph{similarity bias} (viés de similaridade)
e relatam que \num{40} de \num{41} estudos revisados sobre uma mesma base pública de
rolamentos usaram delineamento sujeito a ele.

%% PENDENTE (modelo iCEV): motivo da escolha do tema e trajetória do autor. Texto a ser
%% escrito com informações do autor; entra aqui, antes do problema de pesquisa.

\section{Problema de pesquisa e hipóteses}
\label{sec:intro-problema}

O problema de pesquisa deste trabalho é: quais modos de falha de motores uma cadeia de
medição de baixo custo, formada por acelerômetro MEMS de consumo amostrado a
\SI{1}{\kilo\hertz}, microcontrolador ESP32 e enlace de rádio de longo alcance (LoRa),
permite detectar, e quanto as decisões de estimação alteram o desempenho reportado sobre
essa cadeia? O objeto de estudo é o Kaelix, um dispositivo com esses componentes, que
calcula indicadores de vibração no próprio microcontrolador e envia o valor eficaz e o estado
da máquina por LoRa a um receptor.

A hipótese geral é que uma cadeia desse tipo permite detectar desbalanceamento sob uma taxa
de falso alarme fixada, mas não detectar defeito de rolamento sem desbalanceamento
sobreposto, e que as decisões de estimação alteram o desempenho reportado. Ela se desdobra em seis hipóteses, cada uma com a observação que a refutaria.

\textbf{H1, banda.} Com amostragem a \SI{1}{\kilo\hertz} e o filtro interno do sensor, a
        cadeia preserva as harmônicas da rotação, que caracterizam o desbalanceamento, e
        perde a banda de ressonância em que aparece a assinatura de defeito de rolamento.
        Sob a mesma taxa de falso alarme, a detecção de desbalanceamento deve superar a de
        defeito de rolamento sem massa de desbalanceamento sobreposta. A hipótese é
        refutada se a detecção de rolamento for igual ou maior.

\textbf{H2, integração.} A integração no domínio da frequência, restrita à banda exigida
        pela norma, recupera a velocidade eficaz com erro menor que a integração no
        domínio do tempo, com e sem viés contido na especificação do sensor. A hipótese é
        refutada se a integração no tempo tiver erro igual ou menor em alguma dessas
        condições.

\textbf{H3, limiar.} Com treino apenas em operação normal, o limiar padrão do Isolation
        Forest marca como anomalia uma fração da operação normal maior que a marcada por
        um limiar calibrado, sem alterar a ordem dos escores. A hipótese é refutada se
        essa fração for igual ou menor, ou se a ordem mudar.

\textbf{H4, métrica.} Entre a cadeia de referência e a decimada a \SI{1}{\kilo\hertz}, a área parcial sob
        a curva característica de operação do receptor (ROC), restrita à região de falso
        alarme baixo, cai proporcionalmente mais que a área global. A hipótese é refutada
        se a queda relativa da área parcial for igual ou menor.

\textbf{H5, protocolo.} O mesmo modelo, sobre os mesmos dados e descritores, atinge
        acurácia maior sob partição aleatória por janela que sob partição por ensaio. A
        hipótese é refutada se a acurácia for igual ou menor.

\textbf{H6, dispositivo.} \emandamento{O dispositivo montado adquire janelas a
        \SI{1}{\kilo\hertz} e entrega as leituras ao receptor por LoRa com integridade
        verificada. A hipótese é refutada se as janelas forem descartadas pelo limite de
        \SI{2}{\percent} de desvio de taxa do \emph{firmware} ou se o receptor não validar
        os pacotes pela verificação cíclica de redundância (CRC).}

As taxas de detecção por modo de falha e a comparação entre limiares vêm de dado real: gravações
de uma bancada instrumentada, adquiridas a \SI{50}{\kilo\hertz} e decimadas em \emph{software}
para a taxa do dispositivo. Os resultados sobre banda, integração, métrica e protocolo vêm de sinal
sintetizado e não indicam o desempenho em uma máquina real. \emandamento{Os registros do dispositivo em
bancada cobrem aquisição e transmissão e não incluem motor nem falha induzida. As novas
medições com o dispositivo montado ainda estão em andamento.}

\section{Objetivos}
\label{sec:intro-objetivos}

O objetivo geral é determinar quais modos de falha a cadeia de medição do Kaelix permite
detectar e calibrar o limiar do detector para uma taxa de falso alarme fixada.

Os objetivos específicos são:

\begin{alineas}
    \item quantificar a perda de contraste do indicador de defeito de rolamento entre a
          cadeia de referência e a decimada a \SI{1}{\kilo\hertz}, e verificar se o delineamento
          permite atribuir essa perda à banda ou ao método de demodulação;

    \item medir a detecção por modo de falha em dado real decimado para a taxa do
          dispositivo, sob taxa de falso alarme fixada;

    \item determinar o método de integração de aceleração para velocidade compatível com a
          norma, e quantificar o erro dos métodos alternativos contra um sinal cuja
          integral é conhecida analiticamente~\cite{brandt2014,iso208163_2022};

    \item medir a taxa de falso alarme produzida pelo limiar padrão do Isolation Forest e
          por um limiar calibrado, e o efeito do parâmetro de contaminação sobre ela;

    \item comparar a queda de desempenho medida pela área sob a curva ROC global e pela
          área parcial restrita à região de falso alarme baixo;

    \item comparar o desempenho de um classificador sob partição aleatória por janela e sob
          partição por ensaio;

    \item \emandamento{verificar, no dispositivo montado, a aquisição a
          \SI{1}{\kilo\hertz} e a entrega das leituras por LoRa ao receptor.}
\end{alineas}

\section{Justificativa}
\label{sec:intro-justificativa}

A maior parte dos motores fica sem monitoramento por causa do custo. Na manutenção corretiva, a falha ocorre em momento não escolhido e pode danificar componentes
íntegros. Na preventiva, componentes em bom estado são trocados sem necessidade, e falhas
anteriores ao intervalo continuam sem detecção \cite{mobley2002}. A manutenção preditiva
decide a intervenção pelo estado medido do equipamento, mas o custo de instrumentação a
limita às máquinas críticas. Um nó de baixo custo pode atender os demais motores, e quem
for instalá-lo precisa saber quais falhas ele detecta e com que taxa de falso alarme.

Os resultados publicados são difíceis de comparar. Os trabalhos que propõem esses nós
raramente informam a banda e o piso de ruído do sensor, e \num{40} de \num{41} estudos
revisados por \citeonline{varejao2025} usaram delineamento sujeito a vazamento entre treino
e teste. Nessas condições, o leitor não sabe se o número mede a falha, a montagem ou o
protocolo. Este trabalho mede o efeito dessas decisões. Quem precisa detectar defeito de
rolamento pode usar os resultados para escolher o sensor.

Boa parte das decisões que alteram o desempenho medido é tomada no \emph{software}, o que
liga o tema à Engenharia de Software: o método de integração, o limiar, o
protocolo de validação e a correspondência entre o modelo treinado e o executado no
dispositivo. No Kaelix, as funções de processamento são separadas do acesso ao
\emph{hardware} e testadas no computador, sem o dispositivo. Um teste compara o escore do
detector em C++ com o da biblioteca usada no treino, e os erros são propagados até o
pacote enviado ao receptor (Apêndice~\ref{ape:praticas}).

\section{Classificação da pesquisa}
\label{sec:intro-classificacao}

Segundo os critérios de \citeonline{prodanov2013}, esta pesquisa é aplicada quanto à
natureza, explicativa quanto aos objetivos, quantitativa quanto à abordagem e
bibliográfica e experimental quanto aos procedimentos técnicos. Ela é aplicada porque gera
conhecimento dirigido à solução de um problema específico \cite[p.~51]{prodanov2013}, o de
saber o que uma cadeia de baixo custo permite medir em motores. É explicativa porque
procura identificar os fatores que determinam o fenômeno estudado
\cite[p.~53]{prodanov2013}, neste caso o desempenho reportado, em função da banda, do
método de integração, do limiar, da métrica e do protocolo de validação. A abordagem é
quantitativa porque os resultados são expressos em números e analisados com recursos
estatísticos como percentagem, média e desvio-padrão \cite[p.~69]{prodanov2013}.

Os procedimentos técnicos são dois. A revisão da literatura é pesquisa bibliográfica
\cite[p.~54]{prodanov2013}. As comparações de integração, limiar, métrica e protocolo são
pesquisa experimental, porque manipulam uma variável e mantêm as demais sob controle para
observar seu efeito \cite[p.~57]{prodanov2013}: a mesma senoide passa por cinco variantes
de integração, e o mesmo modelo, sobre os mesmos dados, passa por duas formas de partição.
Esses experimentos são computacionais, sobre sinal sintetizado e dado público de bancada.
\emandamento{As medições com o dispositivo montado, ainda em andamento, acrescentam o
experimento físico.}

\section{Etapas do trabalho}
\label{sec:intro-etapas}

A integração foi testada contra um resultado analítico, de modo que um erro de implementação
aparece como desvio. Os demais experimentos compararam cadeias ou métodos sobre os mesmos dados,
com a classe de cada janela conhecida.

O trabalho teve quatro etapas. A primeira foi a revisão da literatura, com o registro das
lacunas encontradas. A segunda foi o estudo da aquisição: um sinal sintetizado passou pela
cadeia de referência e pelas cadeias decimadas a \SI{1}{\kilo\hertz}, e os descritores foram
comparados. Em paralelo, gravações de bancada adquiridas em alta taxa foram decimadas para
medir a detecção por modo de falha à taxa do dispositivo. A terceira foi o estudo da
estimação, com a varredura do parâmetro de contaminação, a avaliação no ponto de operação
de baixo falso alarme e a comparação entre as duas formas de partição. A quarta foi o
desenvolvimento do \emph{firmware} e do receptor e o ensaio do dispositivo em bancada.

Os dados intermediários e o código que gera os resultados estão no repositório do projeto.

\section{Organização do trabalho}
\label{sec:intro-organizacao}

O restante do trabalho está organizado em quatro seções. A
seção~\ref{cap:fundamentacao} apresenta os conceitos usados no texto e os trabalhos
relacionados: vibração de máquinas rotativas e assinaturas de falha, aquisição de sinais,
avaliação pela norma, detecção de anomalias e sua avaliação, nós sensores e enlace LoRa, e
termina com as lacunas encontradas na literatura. A
seção~\ref{cap:desenvolvimento} descreve os experimentos, o dispositivo, o \emph{firmware}
e o receptor. A seção~\ref{cap:resultados} apresenta os resultados e os compara com a
literatura. A seção~\ref{cap:conclusoes} avalia as hipóteses, indica os limites dos
resultados e lista o trabalho futuro, que inclui o projeto de \emph{hardware} próprio.

%%%%%%%%%%%%%%%%%%%%%%%%%%%%%%%%%%%%%%%%%%%%%%%%%%%%%%%%%%%%%%%%%%%%%%%%%%%%%%%
%% Seção 2 --- Fundamentação teórica e trabalhos relacionados
%%
%% Junta, por decisão do usuário, a fundamentação e a revisão da literatura, antes
%% separadas: cada tema traz o conceito e o que outros autores publicaram sobre ele.
%% O projeto de hardware (invólucro, térmica, bateria) é trabalho futuro e não entra.
%%
%% Os valores do dispositivo vêm de:
%%   rolamento e rotação   docs/relatorios/relatorio-kaelix.tex, tab:parametros
%%   enlace LoRa           src/comms/lora_config.h (SF9, 125 kHz, 4/5, preâmbulo 8)
%%   carga útil            src/comms/lora.h (21 bytes)
%%   MPU6050               folha de dados rev. 3.1 e mapa de registradores
%%                         rev. 4.0, conferidos no documento primário
%%%%%%%%%%%%%%%%%%%%%%%%%%%%%%%%%%%%%%%%%%%%%%%%%%%%%%%%%%%%%%%%%%%%%%%%%%%%%%%

\makeatletter
\@ifundefined{toprule}{%
  \newcommand{\toprule}{\hline}%
  \newcommand{\midrule}{\hline}%
  \newcommand{\bottomrule}{\hline}%
}{}
\@ifundefined{emandamento}{%
  \@ifundefined{textcolor}%
    {\newcommand{\emandamento}[1]{#1}}%
    {\newcommand{\emandamento}[1]{\textcolor{red}{#1}}}%
}{}
\makeatother

\chapter{Fundamentação teórica e trabalhos relacionados}
\label{cap:fundamentacao}

Esta seção apresenta os conceitos usados no restante do texto e o que outros autores publicaram
sobre cada tema, na ordem em que a informação percorre o dispositivo: a vibração da máquina e as
assinaturas das falhas, a aquisição do sinal, a avaliação pela norma, o detector de anomalias e o
nó sensor com seu enlace de rádio. Sempre que possível, os conceitos são aplicados ao Kaelix, para
que os números das seções seguintes possam ser refeitos. A seção termina com as lacunas encontradas
na literatura.

A revisão partiu das escolhas de projeto e dos resultados obtidos e procurou, para cada um, trabalhos
que os confirmassem, os contradissessem ou não tratassem deles. \emandamento{Informar aqui as bases
consultadas, os termos de busca e o período da busca.} Os autores, os títulos e os identificadores
foram conferidos no registro do editor. Quando o texto integral não pôde ser acessado, isso consta em
nota na própria referência. \emandamento{Antes da versão final, conferir o texto das obras citadas e
as páginas de Prodanov e Freitas.}

%%%%%%%%%%%%%%%%%%%%%%%%%%%%%%%%%%%%%%%%%%%%%%%%%%%%%%%%%%%%%%%%%%%%%%%%%%%%%%%
\section{Manutenção de máquinas rotativas}
\label{sec:fund-manutencao}

A manutenção de equipamentos industriais é usualmente classificada em três estratégias, conforme o
momento da intervenção \cite{mobley2002}. Na manutenção \emph{corretiva}, o equipamento opera até
falhar e só então é reparado. Não há custo de planejamento, mas a falha ocorre em momento não
escolhido e pode danificar componentes íntegros. Na manutenção \emph{preventiva}, a intervenção
ocorre em intervalos fixos de tempo ou de uso, definidos pela vida média esperada dos componentes. A
falha inesperada fica menos frequente, mas componentes em bom estado são trocados sem necessidade, e
falhas anteriores ao intervalo continuam sem detecção. Na manutenção \emph{preditiva}, ou baseada em
condição, a intervenção é decidida pelo estado medido do equipamento: acompanha-se uma grandeza
física que reflete sua condição e intervém-se quando ela indica degradação.

\citeonline{jardine2006} descrevem um programa de manutenção baseada em condição em três etapas:
aquisição de dados, processamento dos dados e tomada de decisão. Na aquisição, escolhem-se a grandeza,
o sensor, o ponto de medição e a taxa. No processamento, extraem-se indicadores dos dados brutos. Na
decisão, esses indicadores definem se a máquina precisa de intervenção.

Entre as grandezas usadas no monitoramento de máquinas rotativas, a vibração é a mais difundida,
porque a maior parte dos defeitos mecânicos (desbalanceamento, desalinhamento, folga, desgaste de
rolamentos e de engrenagens) altera as forças dinâmicas no interior da máquina, e essas forças se
transmitem à carcaça, onde podem ser medidas sem interromper a operação \cite{randall2011livro}.

%%%%%%%%%%%%%%%%%%%%%%%%%%%%%%%%%%%%%%%%%%%%%%%%%%%%%%%%%%%%%%%%%%%%%%%%%%%%%%%
\section{Vibração e assinaturas de falha}
\label{sec:fund-vibracao}

A vibração de uma máquina é descrita por grandezas ligadas por derivação e resumida em descritores
escalares. Cada tipo de defeito altera essa vibração de um modo característico, e essas assinaturas
dividem-se em dois grupos, conforme a faixa de frequência em que aparecem.

\subsection{Deslocamento, velocidade e aceleração}

Vibração é o movimento oscilatório de um ponto em torno de uma posição de equilíbrio. O movimento pode
ser descrito por três grandezas relacionadas: o deslocamento $x(t)$, a velocidade
$v(t) = \mathrm{d}x/\mathrm{d}t$ e a aceleração $a(t) = \mathrm{d}v/\mathrm{d}t$. Para um movimento
senoidal de frequência $f$ e amplitude de deslocamento $X$, com $\omega = 2\pi f$,

\begin{equation}
x(t) = X \sin(\omega t), \qquad
v(t) = \omega X \cos(\omega t), \qquad
a(t) = -\omega^2 X \sin(\omega t).
\label{eq:fund-senoide}
\end{equation}

As amplitudes das três grandezas diferem, portanto, por fatores de $\omega$. Para um mesmo deslocamento,
a velocidade cresce proporcionalmente à frequência, e a aceleração, ao quadrado dela. Por isso, a grandeza
escolhida muda o peso de cada frequência: o deslocamento realça as baixas frequências, a aceleração realça
as altas, e a velocidade fica entre as duas. Conhecida uma das grandezas, as outras podem ser obtidas por
derivação ou integração. Um acelerômetro, que mede $a(t)$, pode fornecer a velocidade por integração,
operação discutida na seção~\ref{sec:fund-integracao}.

\subsection{Descritores escalares de um sinal}
\label{sec:fund-descritores}

Um sinal de vibração registrado durante um intervalo é frequentemente resumido por descritores escalares.
Para uma janela de $N$ amostras $x_n$, com média $\bar{x}$ e desvio-padrão $\sigma$, os mais usados são o
valor eficaz, o pico, o pico a pico, o fator de crista (FC) e a curtose \cite{randall2011livro}:

\begin{equation}
x_{\mathrm{rms}} = \sqrt{\frac{1}{N}\sum_{n=1}^{N} x_n^2},
\qquad
x_{\mathrm{pico}} = \max_n |x_n|,
\qquad
\mathrm{FC} = \frac{x_{\mathrm{pico}}}{x_{\mathrm{rms}}},
\qquad
\kappa = \frac{\frac{1}{N}\sum_{n=1}^{N} (x_n-\bar{x})^4}{\sigma^4}.
\label{eq:fund-descritores}
\end{equation}

O valor eficaz ($x_{\mathrm{rms}}$) mede a energia média do sinal e é a base da avaliação de severidade
(seção~\ref{sec:fund-norma}). O valor de pico e o pico a pico (diferença entre o maior e o menor valor da
janela) medem a excursão máxima. O fator de crista e a curtose ($\kappa$) medem a \emph{impulsividade}:
quanto o sinal concentra sua energia em eventos breves e intensos. Uma senoide pura tem
$\mathrm{FC} = \sqrt{2} \approx \num{1,41}$ e $\kappa = \num{1,5}$, e um ruído gaussiano tem $\kappa = 3$.
O \emph{firmware} usa a curtose em excesso, $\kappa - 3$, que vale 0 para esse ruído. Um sinal composto por
impactos espaçados, como o de um defeito localizado em rolamento (seção~\ref{sec:fund-rolamento}), tem fator
de crista e curtose bem maiores, porque poucos valores extremos dominam a média de quarta potência. Esses
descritores não indicam em que frequência a energia está nem a relação entre componentes de frequências
diferentes. Essa informação vem da análise espectral (seção~\ref{sec:fund-espectro}).

\subsection{Defeitos síncronos com a rotação}

Desbalanceamento e desalinhamento produzem forças que se repetem a cada volta do eixo e, por isso, aparecem
em múltiplos inteiros da frequência de rotação $f_r$, as \emph{harmônicas} da rotação
\cite{randall2011livro}. No desbalanceamento, o centro de massa do rotor não coincide com o eixo de rotação,
e a força centrífuga resultante gira com o eixo. A assinatura dominante é a componente em $1\times f_r$, cuja
amplitude cresce com a massa desbalanceada. No desalinhamento, os eixos acoplados não são colineares, e o
acoplamento impõe esforços que variam mais de uma vez por volta. A assinatura típica é o aumento da componente
em $2\times f_r$ em relação à de $1\times f_r$.

Como as rotações industriais usuais estão entre algumas centenas e alguns milhares de rotações por minuto,
essas componentes ficam entre poucos hertz e algumas centenas de hertz. Um motor a \SI{1750}{r/min}, por
exemplo, tem $f_r = \SI{29,17}{\hertz}$, e suas três primeiras harmônicas estão abaixo de \SI{100}{\hertz}.

\subsection{Defeitos em rolamentos}
\label{sec:fund-rolamento}

O rolamento de esferas é composto por uma pista interna, solidária ao eixo, uma pista externa, fixa à carcaça,
e um conjunto de esferas mantidas equidistantes por uma gaiola. Um defeito localizado, como uma trinca ou um
lascamento em uma das pistas, é atingido periodicamente pelas esferas, e cada passagem produz um impacto. A taxa
de repetição desses impactos depende apenas da geometria do rolamento e da rotação. Ela é chamada frequência de
passagem de esferas, na pista externa (BPFO) ou na pista interna (BPFI), conforme a posição do defeito, e vale
\cite{randall2011}

\begin{equation}
\mathrm{BPFO} = \frac{N_b}{2}\, f_r \left(1 - \frac{d}{D_p}\cos\theta\right),
\qquad
\mathrm{BPFI} = \frac{N_b}{2}\, f_r \left(1 + \frac{d}{D_p}\cos\theta\right),
\label{eq:fund-bpfo}
\end{equation}

\noindent
com $N_b$ o número de esferas, $d$ o diâmetro da esfera, $D_p$ o diâmetro primitivo e $\theta$ o ângulo de
contato. Para o rolamento SKF~6205 ($N_b = 9$, $d/D_p = \num{0,203}$, $\theta = 0$) a \SI{1750}{r/min}, as
expressões dão $\mathrm{BPFO} = \SI{104,6}{\hertz}$ e $\mathrm{BPFI} = \SI{157,9}{\hertz}$.

A maior parte da energia do defeito, porém, fica fora da frequência de repetição. Cada impacto é um evento breve,
que excita as ressonâncias da estrutura entre o rolamento e o ponto de medição, tipicamente na faixa de alguns
quilohertz. O sinal resultante é uma sequência de oscilações amortecidas na frequência de ressonância, que se
repetem na taxa BPFO ou BPFI. No espectro, a energia do defeito fica em torno da ressonância, e a frequência
característica aparece como \emph{modulação} dessa portadora de alta frequência \cite{randall2011}. Nos estágios
iniciais do defeito, as linhas na própria BPFO tendem a ser fracas diante das demais componentes da máquina.

Por isso, o procedimento usual de diagnóstico de rolamentos é a \emph{análise de envelope}. O sinal é filtrado em
uma banda em torno da ressonância, calcula-se seu envelope, isto é, a amplitude instantânea, pelo módulo do sinal
analítico obtido com a transformada de Hilbert, e o espectro desse envelope mostra linhas na frequência de repetição
dos impactos \cite{randall2011}. A curtose espectral, isto é, a curtose calculada em função da frequência, é a
ferramenta usada para localizar a banda em que o sinal é mais impulsivo \cite{antoni2007,antoni2006}. Esse
procedimento só funciona se a banda de ressonância tiver sido adquirida.

%%%%%%%%%%%%%%%%%%%%%%%%%%%%%%%%%%%%%%%%%%%%%%%%%%%%%%%%%%%%%%%%%%%%%%%%%%%%%%%
\section{Aquisição do sinal}
\label{sec:fund-sinais}

A taxa de amostragem fixa a maior frequência representável, a duração da janela fixa a resolução espectral e o
transdutor determina a banda e o piso de ruído da medição.

\subsection{Amostragem, frequência de Nyquist e rebatimento}
\label{sec:fund-amostragem}

Um sistema digital registra amostras $a_n = a(n/f_s)$ do sinal contínuo $a(t)$, tomadas à taxa de amostragem
$f_s$. O teorema da amostragem estabelece que um sinal pode ser reconstruído com base em suas amostras se não
contiver componentes acima de $f_s/2$, a \emph{frequência de Nyquist} \cite{brandt2011}. Uma componente acima
desse limite não desaparece na amostragem: suas amostras são idênticas às de uma componente de frequência mais
baixa, dentro da faixa $[0, f_s/2]$. O fenômeno chama-se \emph{rebatimento} (\emph{aliasing}), e a frequência
aparente de uma componente em $f$ é

\begin{equation}
f_{\mathrm{ap}} = \left| f - f_s \cdot \mathrm{round}\!\left(\frac{f}{f_s}\right) \right|.
\label{eq:fund-alias}
\end{equation}

Depois da amostragem, a componente rebatida não pode mais ser separada de uma componente real naquela frequência.
Para evitar isso, é preciso remover o conteúdo acima de $f_s/2$ \emph{antes} da amostragem, com um filtro
passa-baixas chamado filtro \emph{antirrebatimento} (\emph{anti-aliasing}). Pela mesma razão, reduzir a taxa de um
sinal já digitalizado, operação chamada \emph{decimação}, exige filtrar o sinal abaixo da nova frequência de Nyquist
antes de descartar amostras \cite{brandt2011}.

No Kaelix, com $f_s = \SI{1}{\kilo\hertz}$, a frequência de Nyquist é \SI{500}{\hertz}. Uma ressonância de rolamento
em \SI{3,5}{\kilo\hertz} que chegasse sem filtro ao conversor seria observada, pela Equação~\eqref{eq:fund-alias}, em
$|3500 - 1000 \cdot 4| = \SI{500}{\hertz}$, onde não há nenhuma componente física.

\subsection{Análise espectral}
\label{sec:fund-espectro}

A representação de um sinal em função da frequência é obtida pela transformada discreta de Fourier (DFT), calculada
na prática pelo algoritmo da transformada rápida (FFT). Para uma janela de $N$ amostras,

\begin{equation}
X_k = \sum_{n=0}^{N-1} x_n\, e^{-\mathrm{j}2\pi k n / N}, \qquad k = 0, 1, \dots, N-1,
\label{eq:fund-dft}
\end{equation}

\noindent
e o coeficiente $X_k$ corresponde à frequência $k\,f_s/N$ \cite{brandt2011}. A \emph{resolução} em frequência é
$\Delta f = f_s/N$, o inverso da duração da janela, e duas componentes mais próximas que $\Delta f$ não são separadas.
A janela finita também provoca \emph{vazamento espectral}: uma componente cuja frequência não coincide com um múltiplo
de $\Delta f$ espalha energia pelos coeficientes vizinhos. Esse efeito diminui quando o sinal é multiplicado por uma
função de janela que o leve suavemente a zero nas bordas, como a de Hann \cite{brandt2011}. Para sinais aleatórios, o
espectro de uma única janela é uma estimativa muito ruidosa, e usa-se a \emph{densidade espectral de potência}, obtida
pela média dos espectros de várias janelas.

No Kaelix, com $f_s = \SI{1}{\kilo\hertz}$ e janelas de \num{512} amostras, a janela dura \SI{0,512}{\second} e a
resolução é de \SI{1,95}{\hertz}, suficiente para separar as harmônicas de uma rotação de \SI{29}{\hertz}.

\subsection{Acelerômetros piezoelétricos e MEMS}

O transdutor de referência no monitoramento de máquinas é o acelerômetro piezoelétrico: uma massa sísmica apoiada sobre
um cristal que gera carga elétrica proporcional à força e, portanto, à aceleração. Sua frequência de ressonância interna
é elevada, o que lhe dá resposta plana até vários quilohertz, e seu ruído é baixo. Em contrapartida, ele exige
condicionamento de sinal dedicado e custa caro \cite{albarbar2008}.

Os acelerômetros MEMS capacitivos são fabricados pelos mesmos processos dos circuitos integrados. Uma microestrutura
suspensa por molas se desloca sob aceleração, e o deslocamento altera a capacitância entre eletrodos. O conversor
analógico-digital fica no mesmo encapsulamento, e o sensor entrega amostras digitais por barramento serial. O custo é de
uma a duas ordens de grandeza menor que o de um piezoelétrico, mas a banda e o ruído dependem muito do modelo. Há MEMS
projetados para monitoramento de condição, com banda de dezenas de quilohertz, e MEMS de consumo, projetados para orientação
e movimento humano, com banda de algumas centenas de hertz \cite{albarbar2008,fidali2024}. Para decidir se um MEMS serve
para uma aplicação de vibração, consideram-se a banda útil, a densidade espectral de ruído, a faixa de medição e a
estabilidade do zero.

A densidade espectral de ruído, expressa em $\si{\micro g}/\sqrt{\si{\hertz}}$, indica o ruído por unidade de banda. Para
ruído branco, o ruído eficaz na banda $B$ é a densidade multiplicada por $\sqrt{B}$. Quanto mais larga a banda analisada,
maior o ruído total e maior a menor amplitude de vibração que o sensor consegue distinguir.

\subsection{O acelerômetro MPU6050}
\label{sec:fund-mpu6050}

O Kaelix foi planejado para o MPU6050, uma unidade de medição inercial de consumo que combina acelerômetro e giroscópio de
três eixos. As características relevantes para medição de vibração, conforme a folha de dados e o mapa de registradores do
fabricante, estão na Tabela~\ref{tab:fund-mpu6050}.

\begin{table}[!htb]
\IBGEtab{%
\caption{Características do acelerômetro do MPU6050 relevantes para medição de vibração}%
\label{tab:fund-mpu6050}%
}{%
\begin{tabular}{@{}ll@{}}
\toprule
Característica & Valor \\
\midrule
Faixa de medição                         & $\pm2$, $\pm4$, $\pm8$ ou $\pm16$~g, programável \\
Resolução do conversor                   & \num{16} bits \\
Taxa de saída de dados                   & \SIrange{4}{1000}{\hertz}, programável \\
Filtro passa-baixas digital (DLPF)       & \SIrange{5}{260}{\hertz}, programável \\
Densidade espectral de ruído             & \SI{400}{\micro g}/$\sqrt{\si{\hertz}}$ a \SI{10}{\hertz} \\
Tolerância de zero (eixos X e Y)         & $\pm\SI{50}{\milli g}$ \\
\bottomrule
\end{tabular}%
}{%
\fonte{InvenSense (\citeyear{invensense2011ps}, \citeyear{invensense2012rm}).}%
}
\end{table}

Dois pontos da tabela afetam o uso do sensor em vibração. O primeiro é a banda. A taxa máxima de saída do acelerômetro é de
\SI{1}{\kilo\hertz}, o que põe a frequência de Nyquist em \SI{500}{\hertz}, mas quem fixa a banda efetiva é o filtro
passa-baixas digital interno, configurado pelo campo \texttt{DLPF\_CFG} (Tabela~\ref{tab:fund-dlpf}). No MPU6050, na configuração de banda
mais larga, \texttt{DLPF\_CFG}~=~0, que é também o valor após a inicialização, a banda do acelerômetro é de \SI{260}{\hertz},
e não há configuração em que o acelerômetro opere sem esse filtro \cite{invensense2012rm}. O \emph{firmware} usa
\texttt{DLPF\_CFG}~=~1, com banda de \SI{184}{\hertz}. O mapa de registradores informa a banda e o atraso de cada configuração,
mas não a ordem do filtro nem sua atenuação acima da banda, de modo que o fabricante não especifica a proteção contra
rebatimento.

\begin{table}[!htb]
\IBGEtab{%
\caption{Banda e atraso do filtro passa-baixas digital do acelerômetro do MPU6050, por configuração do campo \texttt{DLPF\_CFG}, com taxa de saída de \SI{1}{\kilo\hertz}}%
\label{tab:fund-dlpf}%
}{%
\begin{tabular}{@{}ccc@{}}
\toprule
\texttt{DLPF\_CFG} & Banda (\si{\hertz}) & Atraso (\si{\milli\second}) \\
\midrule
0 & 260 & \num{0,0}  \\
1 & 184 & \num{2,0}  \\
2 & 94  & \num{3,0}  \\
3 & 44  & \num{4,9}  \\
4 & 21  & \num{8,5}  \\
5 & 10  & \num{13,8} \\
6 & 5   & \num{19,0} \\
\bottomrule
\end{tabular}%
}{%
\fonte{\citeonline{invensense2012rm}.}%
}
\end{table}

O segundo ponto é o ruído. Tratando a densidade de \SI{400}{\micro g}/$\sqrt{\si{\hertz}}$ como ruído branco, o ruído eficaz na
banda de \SI{260}{\hertz} é de $400 \cdot \sqrt{260} \approx \SI{6,5}{\milli g}$, aproximadamente
\SI{0,063}{\meter\per\second\squared}. Na banda de \SI{184}{\hertz}, o mesmo cálculo dá \SI{5,4}{\milli g}. A tolerância de zero
de $\pm\SI{50}{\milli g}$ significa que o sensor pode apresentar uma componente contínua de até aproximadamente
\SI{0,5}{\meter\per\second\squared} sem estar fora de especificação. A integração no domínio do tempo acumula esse tipo de desvio
(seção~\ref{sec:fund-integracao}).

\subsection{Limites de banda e de ruído na literatura}
\label{sec:rev-aquisicao}

Os estudos de adequação de sensor apontam o limite superior de banda como a variável mais importante, e nenhum testou um dispositivo
próximo da faixa do Kaelix. A comparação de três acelerômetros MEMS capacitivos com um piezoelétrico de referência reprovou o de banda
mais estreita, de \SI{1500}{\hertz}, para monitoramento de condição \cite{albarbar2008}. Os grupos que diagnosticam rolamento com MEMS
escolhem sensores de dezenas de quilohertz \cite{landi2023,sharma2024}, e uma implementação de baixo custo usa dois acelerômetros: um
de banda estreita para desbalanceamento e outro de \SI{40}{\kilo\hertz} para rolamento \cite{jakobsen2024}. Os produtos comerciais
medem, ao menos em um eixo, uma ordem de grandeza acima do limite de uma IMU de consumo
\cite{bentlynevada2025,erbessd2025,tractian2022}.

O piso de ruído também limita a medição, independentemente da banda. A densidade de ruído da folha de dados de uma IMU de consumo,
integrada sobre a banda normativa, fica acima do ruído admissível para separar as zonas mais baixas da norma. Os estudos sobre MEMS
apontam a falta de calibração rastreável desses dispositivos \cite{albarbar2008,rodrigues2021}, e a avaliação de sensibilidade
diagnóstica de sensores digitais de catálogo chega à mesma conclusão \cite{fidali2024}. \emandamento{Não foi possível confirmar em
fonte primária a estimativa de ruído admissível usualmente atribuída à literatura de fabricante \cite{looney2017}, e a comparação com
a densidade medida ainda não foi feita.}

O filtro antirrebatimento é pouco discutido nessa literatura, embora afete dispositivos de consumo. A busca não encontrou, na literatura
revisada por pares sobre diagnóstico de rolamento, descrição experimental do efeito do rebatimento sobre descritores como curtose e fator
de crista. O efeito aparece apenas em notas de aplicação de fabricante, que não foram conferidas nesta revisão.

%%%%%%%%%%%%%%%%%%%%%%%%%%%%%%%%%%%%%%%%%%%%%%%%%%%%%%%%%%%%%%%%%%%%%%%%%%%%%%%
\section{Avaliação da severidade pela norma}
\label{sec:fund-norma}

A avaliação da vibração de máquinas industriais medida em partes não rotativas é normalizada pela ISO 20816-3:2022, que substituiu a
ISO 10816-3:2009 \citemesmoautor{International Organization for Standardization}{iso108163_2009}{iso208163_2022}. A norma aplica-se a
máquinas com potência acima de \SI{15}{\kilo\watt} e rotação entre \SI{120}{r/min} e \SI{30000}{r/min} e adota como grandeza de avaliação
a \emph{velocidade eficaz de banda larga}. Segundo o texto da norma, vibrações com o mesmo valor eficaz de velocidade na faixa de
\SIrange{10}{1000}{\hertz} podem, em geral, ser consideradas de mesma severidade, ao passo que critérios em deslocamento ou em aceleração
variariam com a frequência \cite{iso208163_2022}.

Pela mesma razão, a norma exige do instrumento resposta plana em \emph{ao menos} \SIrange{10}{1000}{\hertz}, com limite inferior de até
\SI{2}{\hertz} para máquinas com rotação próxima ou abaixo de \SI{600}{r/min}, e determina que o sinal de acelerômetros montados na carcaça
seja integrado para fornecer velocidade. Um instrumento com banda menor que a exigida não atende à norma, mesmo que o processamento posterior
esteja correto. A edição anterior delega a definição de resposta plana a uma norma de instrumentação, que usa tolerância de banda passante de
aproximadamente \SI{10}{\percent} entre \SI{12,6}{\hertz} e \SI{1000}{\hertz} e limitação de banda por um par de filtros de Butterworth de
terceira ordem \cite{iso2954_2012}.

O valor medido é comparado com quatro zonas de avaliação:

\begin{alineas}
  \item zona A: vibração típica de máquinas recém-comissionadas;
  \item zona B: vibração aceitável para operação de longo prazo sem restrição;
  \item zona C: vibração insatisfatória para operação contínua de longo prazo, em que a máquina pode operar por período limitado, até que haja
        oportunidade de ação corretiva;
  \item zona D: vibração suficientemente severa para causar dano à máquina.
\end{alineas}

Os limites entre as zonas dependem do porte da máquina e do tipo de fundação, conforme a Tabela~\ref{tab:fund-zonas}. A fundação é dita rígida
quando a primeira frequência natural do conjunto máquina--fundação está suficientemente acima da frequência de excitação principal, e flexível
no caso contrário. Os valores de zona são iguais nas edições de 2009 e de 2022.

\begin{table}[!htb]
\IBGEtab{%
\caption{Limites entre zonas de avaliação para máquinas industriais, em velocidade eficaz de banda larga}%
\label{tab:fund-zonas}%
}{%
\begin{tabular}{@{}llccc@{}}
\toprule
Grupo & Fundação & \multicolumn{3}{c}{Limite entre zonas (\si{\milli\meter\per\second})} \\
\cmidrule{3-5}
 & & A/B & B/C & C/D \\
\midrule
1 & rígida   & \num{2,3} & \num{4,5} & \num{7,1}  \\
1 & flexível & \num{3,5} & \num{7,1} & \num{11,0} \\
2 & rígida   & \num{1,4} & \num{2,8} & \num{4,5}  \\
2 & flexível & \num{2,3} & \num{4,5} & \num{7,1}  \\
\bottomrule
\end{tabular}%
}{%
\fonte{\citeonline{iso208163_2022}.}%
\nota{Grupo 1: potência acima de \SI{300}{\kilo\watt}, ou máquinas elétricas com altura de eixo de pelo menos \SI{315}{\milli\meter}. Grupo 2: potência de \SIrange{15}{300}{\kilo\watt}, ou máquinas elétricas com altura de eixo entre \SI{160}{\milli\meter} e \SI{315}{\milli\meter}.}%
}
\end{table}

\subsection{Integração de aceleração para velocidade}
\label{sec:fund-integracao}

Como a norma avalia a severidade em velocidade e o acelerômetro fornece aceleração, é preciso integrar o sinal. Há duas formas de fazer isso
\cite{brandt2014}.

No \emph{domínio do tempo}, a velocidade é aproximada pela soma acumulada das amostras de aceleração multiplicadas pelo intervalo de amostragem. O
método é simples, mas qualquer componente contínua presente na aceleração, como um desvio de zero do sensor, é integrada como uma rampa, e o erro
cresce linearmente com a duração do registro. Um filtro passa-alta aplicado antes da integração reduz esse efeito, mas introduz transitórios nas
bordas da janela.

No \emph{domínio da frequência}, usa-se a relação da Equação~\eqref{eq:fund-senoide}: cada componente de frequência $f$ da velocidade é a componente
correspondente da aceleração dividida por $\mathrm{j}2\pi f$. Calcula-se a DFT da aceleração, divide-se cada coeficiente por $\mathrm{j}2\pi f_k$,
anulam-se os coeficientes fora da banda de interesse e retorna-se ao tempo pela transformada inversa. Anular esses coeficientes elimina a componente
contínua e a divisão por zero em $f = 0$. \citeonline{brandt2014} compararam os dois métodos e concluíram que a integração no domínio da frequência é
a mais exata para registros longos. A busca não encontrou a aplicação desse método à velocidade eficaz de banda larga em um dispositivo de banda
limitada, ou seja, quanto um instrumento desse tipo subestima o indicador da norma.

Como exemplo, uma aceleração senoidal de amplitude \SI{3,0}{\meter\per\second\squared} a \SI{50}{\hertz} corresponde a uma velocidade de amplitude
$3{,}0/(2\pi \cdot 50) \approx \SI{9,55}{\milli\meter\per\second}$ e valor eficaz $\num{9,55}/\sqrt{2} \approx \SI{6,75}{\milli\meter\per\second}$. Pela
Tabela~\ref{tab:fund-zonas}, esse valor está na zona B para máquinas do grupo 1 em fundação flexível, na zona C para o grupo 1 rígido e o grupo 2
flexível, e na zona D para o grupo 2 rígido. Um erro de poucos milímetros por segundo muda, portanto, a classificação da máquina. A
seção~\ref{cap:resultados} usa esse valor analítico para comparar os métodos de integração.

%%%%%%%%%%%%%%%%%%%%%%%%%%%%%%%%%%%%%%%%%%%%%%%%%%%%%%%%%%%%%%%%%%%%%%%%%%%%%%%
\section{Detecção de anomalias e sua avaliação}
\label{sec:fund-anomalias}

Um detector de anomalias depende de três escolhas: o paradigma de aprendizado, o algoritmo que produz o escore e o limiar que transforma o escore em
decisão. O resultado informado também depende da métrica e da separação entre dados de treino e de teste.

\subsection{Paradigmas de aprendizado}

A condição da máquina pode ser avaliada de modos diferentes, conforme os dados disponíveis para treinar o modelo \cite{chandola2009}. Na abordagem
\emph{supervisionada}, o modelo aprende com exemplos rotulados de cada condição (normal e cada tipo de falha) e classifica novos dados em uma dessas
classes. Na abordagem \emph{semissupervisionada}, ou de classe única, o modelo aprende apenas com exemplos da condição normal e marca como anomalia o que
se afasta dela. Na abordagem \emph{não supervisionada}, não há rótulo algum, e o modelo supõe que as anomalias são raras e diferentes da maioria dos dados.

Em manutenção preditiva, a abordagem supervisionada esbarra em um problema prático: uma máquina recém-instalada não tem histórico de falhas, e reunir
exemplos rotulados de cada defeito exige que eles ocorram. Os dados de operação normal existem desde o primeiro dia. Por isso, um dispositivo instalado em
máquina sem histórico precisa treinar só com eles \cite{chandola2009}. Em contrapartida, o modelo detecta
\emph{afastamento} da normalidade e não identifica \emph{qual} defeito está presente.

\subsection{Isolation Forest}
\label{sec:fund-iforest}

O Isolation Forest \cite{liu2008} se baseia na ideia de que anomalias são poucas e diferentes e, por isso, mais fáceis de \emph{isolar} do resto dos dados do
que os pontos normais. O algoritmo constrói um conjunto de árvores de isolamento. Cada árvore usa uma subamostra aleatória dos dados de treino e a divide
recursivamente: em cada nó, escolhe-se ao acaso um atributo e um valor de corte entre o mínimo e o máximo desse atributo, até que cada ponto fique isolado
em uma folha ou que se atinja um limite de profundidade. Um ponto anômalo, por estar em região de baixa densidade, tende a ser isolado em poucas divisões,
perto da raiz. Um ponto normal exige mais divisões.

O comprimento do caminho $h(x)$ entre a raiz e a folha em que o ponto $x$ cai mede, portanto, quão fácil é isolá-lo. Para comparar árvores construídas com
subamostras de tamanho $\psi$, essa medida é normalizada pelo comprimento médio esperado do caminho em uma árvore binária de busca com $\psi$ pontos,

\begin{equation}
c(\psi) = 2H(\psi-1) - \frac{2(\psi-1)}{\psi}, \qquad H(i) \approx \ln i + \num{0,5772},
\label{eq:fund-cn}
\end{equation}

\noindent
e o escore de anomalia é

\begin{equation}
s(x, \psi) = 2^{-\,\mathbb{E}[h(x)] / c(\psi)},
\label{eq:fund-score}
\end{equation}

\noindent
em que $\mathbb{E}[h(x)]$ é a média do comprimento do caminho sobre todas as árvores \cite{liu2008}. O escore está entre 0 e 1. Valores próximos de 1 indicam
ponto fácil de isolar, candidato a anomalia. Valores bem abaixo de \num{0,5} indicam ponto normal. Se todos os pontos têm escore próximo de \num{0,5}, os dados
não contêm anomalias distintas. Os autores recomendam $t = 100$ árvores e subamostras de $\psi = 256$ pontos, para as quais $c(256) \approx \num{10,2}$.

Algumas características favorecem a execução embarcada: o treino usa subamostras pequenas, a inferência percorre cada árvore uma única vez com uma comparação
por nó, e o modelo treinado é inteiramente descrito por atributos e valores de corte \cite{liu2008}.

\subsection{Escolha do detector}
\label{sec:rev-detector}

Dois comparativos chegam a conclusões diferentes. O primeiro avalia \num{30} algoritmos de vários paradigmas
sobre \num{57} conjuntos de dados e conclui que \emph{nenhum} método não supervisionado é estatisticamente superior aos demais \cite{han2022}. O segundo compara
\num{33} algoritmos não supervisionados sobre \num{52} conjuntos tabulares reais e encontra vantagem para o Extended Isolation Forest, com a ressalva de que os
dados se separam em dois regimes \cite{bouman2024}.

Este trabalho usa o algoritmo original por causa do custo embarcado. Um nó da variante estendida guarda um vetor normal com tantas componentes quantos forem os
descritores, e a inferência passa de uma comparação para um produto interno por nó. \emandamento{Esse é um argumento de projeto. O consumo de memória e de tempo
ainda não foi medido no dispositivo.}

Para máquinas rotativas, o algoritmo original tem bons resultados publicados. Entre \num{11} detectores não supervisionados, ele atingiu F1 de \SI{97,19}{\percent}
em rolamento \emph{run-to-failure}, de \SI{99,71}{\percent} em caixa de engrenagens e de \SI{97,20}{\percent} em falhas mecânicas de bancada \cite{brito2022}. Os dois
primeiros conjuntos foram adquiridos a \SI{20}{\kilo\hertz}, banda que uma IMU de consumo não tem.

\subsection{Limiar de decisão e taxa de falso alarme}
\label{sec:fund-limiar}

O escore contínuo de anomalia precisa ser convertido em decisão por um \emph{limiar}: acima dele, o ponto é marcado como anômalo. A escolha do limiar troca um tipo de
erro pelo outro. Um limiar baixo detecta mais falhas e marca mais pontos normais como anomalia, e um limiar alto faz o contrário.

Na implementação de referência do algoritmo, a escolha é controlada pelo parâmetro de contaminação, que representa a fração esperada de anomalias nos dados de treino.
No valor padrão, o limiar corresponde à referência de \num{0,5} do escore, sem calibração para os dados. Com um valor numérico, o limiar é posicionado de modo que essa
fração dos pontos de treino fique acima dele. Nos dois casos, o parâmetro muda apenas um deslocamento fixo e não muda o escore de cada ponto \cite{liu2008,scikitlearn2026}.
Por isso, toda métrica de ordenação, inclusive a área sob a curva ROC (seção~\ref{sec:fund-roc}), não depende desse parâmetro.

Quando o treino contém apenas operação normal, o limiar pode ser calibrado diretamente pela taxa de falso alarme desejada: se o limiar é o quantil de ordem $1-\alpha$ dos
escores de dados normais não usados no ajuste das árvores, espera-se que uma fração $\alpha$ da operação normal fique acima dele. Outra alternativa modela apenas a cauda da
distribuição dos escores pela teoria de valores extremos e fixa o limiar com base em um risco alvo, isto é, na probabilidade de falso alarme admitida \cite{siffer2017}.

Os trabalhos que relatam bom desempenho não usam o valor padrão. Um fixa o escore normalizado em \num{0,75} sob a hipótese de treino limpo \cite{antonini2023}, e outro
calcula o limiar pela razão de contaminação conhecida \cite{brito2022}. As duas implementações embarcadas do mesmo detector que informam o falso alarme, ou permitem
calculá-lo, ficam entre \SI{22}{\percent} e \SI{38}{\percent} sobre operação normal \cite{chevtchenko2023,kolok2025}. Com essas taxas, um dispositivo emitiria de um a dois
alarmes falsos a cada cinco avaliações de uma máquina em operação normal.

\subsection{Matriz de confusão e taxas derivadas}

Com a anomalia como classe positiva, cada decisão do detector cai em uma de quatro categorias: verdadeiro positivo (VP), anomalia corretamente detectada; falso positivo (FP),
operação normal marcada como anomalia; verdadeiro negativo (VN); e falso negativo (FN), anomalia não detectada. Dessas contagens derivam a taxa de verdadeiros positivos (TPR),
a taxa de falsos positivos (FPR), a precisão e a medida $F_1$ \cite{fawcett2006}:

\begin{equation}
\mathrm{TPR} = \frac{\mathrm{VP}}{\mathrm{VP} + \mathrm{FN}},
\quad
\mathrm{FPR} = \frac{\mathrm{FP}}{\mathrm{FP} + \mathrm{VN}},
\quad
\text{precisão} = \frac{\mathrm{VP}}{\mathrm{VP} + \mathrm{FP}},
\quad
F_1 = \frac{2\cdot \text{precisão} \cdot \mathrm{TPR}}{\text{precisão} + \mathrm{TPR}}.
\label{eq:fund-taxas}
\end{equation}

A TPR é a \emph{taxa de detecção}, e a FPR é a \emph{taxa de falso alarme}. Em monitoramento contínuo, a taxa de falso alarme pesa na operação: um dispositivo que avalia a máquina a
cada poucos minutos e marca \SI{10}{\percent} da operação normal como anomalia gera dezenas de alarmes falsos por dia, e alarmes frequentes demais deixam de ser atendidos. A acurácia,
fração de decisões corretas, informa pouco quando as classes são desbalanceadas, porque um detector que nunca dispara alcança acurácia alta se as anomalias são raras.

\subsection{Curva ROC, área sob a curva e área parcial}
\label{sec:fund-roc}

Como o limiar é uma escolha, um detector tem uma TPR para cada limiar. A curva característica de operação do receptor (ROC) traça a TPR contra a FPR à medida que o limiar varre todos os
valores possíveis, e a área sob essa curva (AUC) resume o desempenho em um número entre \num{0,5}, o de uma decisão ao acaso, e 1, o de um detector perfeito \cite{fawcett2006}. A AUC é a
probabilidade de que uma anomalia escolhida ao acaso receba escore maior que um ponto normal escolhido ao acaso. Como depende apenas da ordem dos escores, ela não muda quando o limiar é
deslocado.

Um dispositivo de campo, porém, opera em um único ponto da curva, com falso alarme baixo, e dois detectores com a mesma AUC podem ter desempenhos muito diferentes nessa região. A
\emph{área parcial sob a curva} (pAUC) restringe a integração à faixa de FPR de interesse, entre 0 e um limite $\beta$, e a área nessa faixa é
$A_\beta = \int_{0}^{\beta} \mathrm{TPR}\; \mathrm{d}\,\mathrm{FPR}$. Este trabalho usa a forma padronizada de \citeonline{mcclish1989}, também calculada pela biblioteca usada no treino,

\begin{equation}
\mathrm{pAUC}(\beta) = \frac{1}{2}\left(1 + \frac{A_\beta - \beta^2/2}{\beta - \beta^2/2}\right),
\label{eq:fund-pauc}
\end{equation}

\noindent
que vale \num{0,5} para uma decisão ao acaso e 1 para um detector perfeito. O desafio de referência de detecção não supervisionada de anomalias em máquinas usa outra forma, $A_\beta/\beta$, com
$\beta = \num{0,1}$ \cite{koizumi2020}. Outro trabalho informa a área sob a curva de precisão e revocação (PR-AUC) além do F1 \cite{brito2022}.

\subsection{Validação cruzada e vazamento de informação}
\label{sec:fund-validacao}

O desempenho de um modelo precisa ser medido em dados não usados no treino. Na \emph{validação cruzada em $k$ partições}, os dados são divididos em $k$ partições, e o modelo é treinado $k$ vezes,
cada vez com uma partição reservada para teste. O desempenho é a média das $k$ avaliações, e sua dispersão indica quanto o resultado depende dos dados usados no teste.

A validação cruzada exige que os dados de treino e de teste sejam independentes. Quando não são, ocorre \emph{vazamento de informação}: o modelo tem acesso, no treino, a informação que não estaria
disponível na aplicação real, e o desempenho medido fica acima do verdadeiro \cite{kaufman2012}. Em sinais de vibração, a forma mais comum de vazamento vem da segmentação. Um ensaio de alguns
segundos é dividido em muitas janelas, e janelas do mesmo ensaio têm características que nada têm a ver com a falha, como nível de ruído, ganho do sensor e condição de fixação. Se as janelas são
sorteadas entre treino e teste, o modelo pode acertar porque reconhece o ensaio. \citeonline{varejao2025} chamam esse efeito de \emph{viés de similaridade}. A correção é a \emph{validação por
grupos}, aqui chamada de partição por ensaio: todas as janelas de um mesmo ensaio ficam na mesma partição, e o teste é sempre feito sobre ensaios que o modelo não viu.

Esse problema já foi publicado e não é contribuição deste trabalho. De \num{41} estudos revisados sobre uma mesma base pública de rolamentos, \num{40} usaram delineamento sujeito a esse vazamento
\cite{varejao2025}, e a diferença medida é da ordem de dezenas de pontos percentuais \cite{vieira2026}. Outras revisões sobre classificação de falha de rolamento chegam a conclusão parecida
\cite{abburi2023,smith2015}.

%%%%%%%%%%%%%%%%%%%%%%%%%%%%%%%%%%%%%%%%%%%%%%%%%%%%%%%%%%%%%%%%%%%%%%%%%%%%%%%
\section{Nós sensores e enlace LoRa}
\label{sec:fund-no}

O Kaelix é um nó sensor: um sistema embarcado que mede a vibração e envia o resultado por rádio.

\subsection{Sistemas embarcados para monitoramento}
\label{sec:fund-embarcados}

Um nó sensor de monitoramento é um sistema embarcado: um computador dedicado a uma única função, integrado ao equipamento que atende, com pouca memória e pouca capacidade de processamento. Seu
componente central é o microcontrolador, que reúne processador, memória e periféricos (conversores, interfaces seriais, temporizadores) em um único circuito integrado.

O Kaelix usa microcontroladores da família ESP32. O ESP32-S3, por exemplo, tem dois núcleos de até \SI{240}{\mega\hertz}, \SI{512}{\kilo\byte} de memória estática de acesso aleatório (SRAM) e
rádios Wi-Fi e Bluetooth integrados \cite{espressif2026}.

Executar a detecção no próprio dispositivo reduz o uso do rádio, porque transmitir uma janela de \num{512} amostras de três eixos ocuparia o enlace por muito mais tempo que transmitir o resultado. A
execução de modelos de aprendizado de máquina em microcontroladores, com memória de centenas de quilobytes e consumo de miliwatts, é o campo denominado TinyML \cite{antonini2023,barbariol2022}.

\subsection{Nós de baixo custo publicados}
\label{sec:rev-implantacao}

Os nós sem fio para manutenção preditiva dividem-se em dois grupos. Os produtos comerciais têm banda uma ordem de grandeza acima da de uma IMU de consumo
\cite{bentlynevada2025,erbessd2025,tractian2022}. Os nós acadêmicos de baixo custo usam a mesma arquitetura estudada aqui: microcontrolador, IMU MEMS com taxa na ordem de \SI{1}{\kilo\hertz} e
detecção não supervisionada \cite{antonini2023,bisio2025,kolok2025}.

O trabalho mais próximo em \emph{hardware} e em modo de falha relata acurácia de \SI{72,9 \pm 2,4}{\percent}, com validação cruzada de cinco partições sobre cerca de \num{900} amostras e \emph{um
único} modo de falha \cite{kolok2025}. O trabalho termina no valor eficaz de aceleração, sem integração para velocidade e sem critério da norma.

A comparação mais direta é embarcada e mede três detectores no mesmo nó e no mesmo conjunto \cite{chevtchenko2023}. Nela, um método baseado em vizinhança superou o algoritmo usado aqui em
especificidade e em latência, com execução aproximadamente \num{26} vezes mais rápida. \emandamento{O método de vizinhança precisa manter o conjunto de referência na memória, enquanto a floresta
guarda apenas limiares e índices, e esse foi o motivo para manter o detector atual. O consumo de memória no dispositivo ainda não foi medido.}

O mesmo trabalho mostra que um motor com cerca de um ano de uso tem assinatura normal diferente da de motores novos \cite{chevtchenko2023}. Um modelo único com limiar fixo sofre com essa mudança,
tanto entre máquinas quanto ao longo da vida útil.

Outros trabalhos mostram que o detector treina e executa dentro de um microcontrolador dessa classe \cite{antonini2023}, que há variante criada para reduzir seu custo em TinyML \cite{barbariol2022}
e que dispositivos da mesma família suportam processamentos mais pesados \cite{kumar2026,uselis2025}.

\subsection{Enlace LoRa}
\label{sec:fund-lora}

LoRa é uma técnica de modulação de rádio criada para enlaces de longo alcance e baixo consumo, com baixa taxa de dados \cite{augustin2016}. Ela usa espalhamento espectral por \emph{chirp}: cada
símbolo é representado por um sinal cuja frequência varia linearmente ao longo de toda a largura de banda do canal. O espalhamento permite demodular sinais abaixo do nível de ruído, o que aumenta o
alcance. LoRa designa a camada física, e LoRaWAN é um protocolo de rede definido sobre ela, com servidores de rede e gerência de dispositivos. Um enlace ponto a ponto, entre o dispositivo e um
receptor próprio, pode usar a modulação LoRa sem o protocolo LoRaWAN.

Três parâmetros definem o equilíbrio entre alcance, taxa e consumo \cite{semtech2015}:

\begin{alineas}
  \item fator de espalhamento (SF), de 6 a 12: cada símbolo carrega SF bits e é formado por $2^{\mathrm{SF}}$ \emph{chips}, e cada incremento de SF dobra a duração do símbolo e aumenta a
        sensibilidade do receptor;
  \item largura de banda (BW), tipicamente \SI{125}{\kilo\hertz}: banda mais larga encurta o símbolo e reduz a sensibilidade;
  \item taxa de codificação (CR), de 4/5 a 4/8: fração de redundância acrescentada para correção de erros.
\end{alineas}

A taxa de símbolos é $R_s = \mathrm{BW}/2^{\mathrm{SF}}$, e a duração do símbolo é $T_s = 1/R_s$. A duração total de um pacote, o \emph{tempo no ar}, é a soma da duração do preâmbulo e da carga útil
\cite{semtech2015}:

\begin{equation}
\begin{split}
T_{\mathrm{pacote}} = {} & (n_{\mathrm{pre}} + \num{4,25})\,T_s \\
& + \left[ 8 + \max\!\left( \left\lceil \frac{8\,\mathrm{PL} - 4\,\mathrm{SF} + 28 + 16\,\mathrm{CRC} - 20\,\mathrm{IH}}{4(\mathrm{SF} - 2\,\mathrm{DE})} \right\rceil (\mathrm{CR} + 4),\; 0 \right) \right] T_s,
\end{split}
\label{eq:fund-toa}
\end{equation}

\noindent
em que $n_{\mathrm{pre}}$ é o número de símbolos do preâmbulo, PL o número de bytes da carga útil, CRC vale 1 quando a verificação de redundância está habilitada, IH vale 1 quando o cabeçalho é
omitido, DE vale 1 quando a otimização para baixa taxa está ativa, e CR vale de 1 (para 4/5) a 4 (para 4/8).

O enlace do Kaelix usa SF9, BW~=~\SI{125}{\kilo\hertz}, CR~=~4/5, preâmbulo de 8 símbolos e carga útil de \num{21} bytes, com cabeçalho e CRC habilitados e sem otimização para baixa taxa. A duração
do símbolo é $2^9/\num{125000} = \SI{4,096}{\milli\second}$, o preâmbulo ocupa $(8 + \num{4,25}) \times \num{4,096} \approx \SI{50,2}{\milli\second}$, e a carga útil ocupa
$8 + \lceil 176/36 \rceil \times 5 = 33$ símbolos, ou \SI{135,2}{\milli\second}. O tempo no ar resultante é de aproximadamente \SI{185}{\milli\second}. Em SF12, a mesma carga útil ocuparia o rádio
por um tempo várias vezes maior, em troca de mais alcance.

%%%%%%%%%%%%%%%%%%%%%%%%%%%%%%%%%%%%%%%%%%%%%%%%%%%%%%%%%%%%%%%%%%%%%%%%%%%%%%%
\section{Síntese e lacunas}
\label{sec:rev-lacunas}

A literatura revisada tem mais resultados sobre aquisição, em que um sensor de consumo tem maiores limitações, do que sobre estimação, em que este trabalho concentra seus experimentos.

Os trabalhos revisados apoiam quatro pontos usados aqui. O limiar padrão não é calibrado, e há uma alternativa baseada na teoria de valores extremos. O \emph{benchmark} de referência informa a métrica
no ponto de operação. O detector já foi executado em microcontrolador. A norma exige a integração de aceleração para velocidade, e o método correto é conhecido.

Três escolhas deste trabalho vão contra o que a literatura recomenda. Nenhum estudo de adequação revisado testou sensor com limite superior próximo ao usado aqui, e a assinatura de rolamento depende de
uma banda que o dispositivo não adquire. O piso de ruído de uma IMU de consumo limita a medição, qualquer que seja o processamento. Por fim, a maior comparação de métodos não supervisionados recomenda
outra variante do detector, e um método de vizinhança foi mais específico e mais rápido no mesmo tipo de \emph{hardware} embarcado. \emandamento{Este trabalho ainda não mediu o consumo de memória e de
tempo no dispositivo, que é o argumento para manter o detector atual.}

A revisão encontrou sete lacunas. Duas lacunas supostas no início da revisão não se confirmaram: dois trabalhos informam o falso alarme do
detector em nó embarcado, ou permitem calculá-lo \cite{chevtchenko2023,kolok2025}, e, com isso, a fração da operação normal marcada pelo limiar padrão já está quantificada.

\textbf{Lacuna 1: efeito isolado do parâmetro de contaminação.} Não se encontrou trabalho que variasse o parâmetro sobre a \emph{mesma} base e o \emph{mesmo} modelo e informasse a fração da operação normal
        marcada.

\textbf{Lacuna 2: falso alarme de detector treinado no próprio dispositivo.} O único trabalho revisado que treina no próprio dispositivo deixa a avaliação de acurácia fora do escopo \cite{antonini2023}.

\textbf{Lacuna 3: contraste diagnóstico em função da banda.} A literatura discute o efeito da banda de forma qualitativa, e o único estudo revisado que varre a taxa de amostragem mede a acurácia de um
        classificador supervisionado, e não o contraste do indicador \cite{elbouharrouti2024}.

\textbf{Lacuna 4: rebatimento sobre descritores estatísticos.} O efeito não foi encontrado em literatura revisada por pares e aparece apenas em notas de aplicação de fabricante.

\textbf{Lacuna 5: efeito da limitação de banda sobre o indicador da norma.} O erro dos métodos de integração está quantificado \cite{brandt2014}, mas não sua aplicação a esse indicador em dispositivo de banda
        limitada.

\textbf{Lacuna 6: detecção limitada pelo piso de ruído.} Não se encontrou análise de defeito mínimo detectável em função da densidade de ruído para IMUs de consumo.

\textbf{Lacuna 7: separação entre efeito de banda e efeito de montagem.} Não se encontrou \emph{benchmark} aberto com sinais gravados \emph{simultaneamente} por MEMS de banda limitada e por piezoelétrico no
        mesmo ponto.

A seção~\ref{cap:resultados} trata das lacunas 1 e 4 em sinal sintetizado e oferece à lacuna 3 apenas dois pontos, insuficientes para formar uma função da banda. As demais continuam abertas, e parte delas
aparece como trabalho futuro na seção~\ref{sec:conc-futuro}.

%%%%%%%%%%%%%%%%%%%%%%%%%%%%%%%%%%%%%%%%%%%%%%%%%%%%%%%%%%%%%%%%%%%%%%%%%%%%%%%
%% Capítulo 3 --- Desenvolvimento e métodos
%%
%% Descreve os experimentos e o que está em kaelix-firmware: componentes reais,
%% firmware, receptor e ensaio de bancada. O projeto de hardware (placas,
%% invólucro, térmica, bateria) é trabalho futuro, por decisão do usuário.
%%
%% Valores conferidos em platformio.ini, src/sensors/vibration.h,
%% src/comms/lora_config.h, src/comms/lora.h, src/main.cpp e logs/.
%%%%%%%%%%%%%%%%%%%%%%%%%%%%%%%%%%%%%%%%%%%%%%%%%%%%%%%%%%%%%%%%%%%%%%%%%%%%%%%

\makeatletter
\@ifundefined{toprule}{%
  \newcommand{\toprule}{\hline}%
  \newcommand{\midrule}{\hline}%
  \newcommand{\bottomrule}{\hline}%
}{}
\@ifundefined{emandamento}{%
  \@ifundefined{textcolor}%
    {\newcommand{\emandamento}[1]{#1}}%
    {\newcommand{\emandamento}[1]{\textcolor{red}{#1}}}%
}{}
\makeatother

\chapter{Desenvolvimento e métodos}
\label{cap:desenvolvimento}

O trabalho tem duas partes. A primeira são os experimentos computacionais, que usam sinal
sintetizado e gravações públicas de bancada para avaliar a aquisição e a estimação. A
segunda é o dispositivo: um acelerômetro e um rádio LoRa ligados a um microcontrolador
ESP32, o \emph{firmware} que os controla e o receptor.

\section{Delimitação}
\label{sec:dev-delimitacao}

O objeto de estudo é o Kaelix, um dispositivo de monitoramento de vibração de baixo custo.
Este trabalho trata da cadeia de medição, do \emph{firmware} e da comunicação. O projeto de
placas de circuito impresso, o invólucro, a bateria e o comportamento térmico ficam para
trabalho futuro (seção~\ref{sec:conc-futuro}). Nos ensaios, os componentes foram ligados
em montagem de bancada e alimentados pela porta USB.

A norma ISO 20816-3 serve aqui como referência para comparar métodos de integração. Ela
exige resposta plana de \SIrange{10}{1000}{\hertz}, e o Kaelix não atende a esse requisito,
porque o filtro interno do sensor limita a banda a \SI{184}{\hertz}. Por isso, o
dispositivo não classifica máquinas nas zonas da norma.

\section{Modelo de sinal e cadeias de aquisição}
\label{sec:dev-sinal}

O efeito da aquisição foi medido sobre dois tipos de sinal: um sinal sintetizado, cuja
resposta correta é conhecida, e gravações reais decimadas para a taxa do dispositivo.

\subsection{Sinal sintetizado}

O sinal sintetizado combina dois mecanismos. Desbalanceamento e desalinhamento aparecem como
harmônicas da rotação, todas abaixo de \SI{100}{\hertz}. O defeito de pista externa aparece
como um trem de impulsos amortecidos, repetido na taxa característica da geometria, que
modula uma ressonância estrutural. A aceleração simulada é

\begin{equation}
a(t) \;=\; \underbrace{\sum_{k=1}^{3} A_k \sin\!\left(2\pi k f_r t + \varphi_k\right)}_{\text{banda baixa}}
\;+\; \underbrace{\sum_{i} D\, e^{-\zeta_d (t-t_i)} \sin\!\left(2\pi f_{\mathrm{res}} (t-t_i)\right) u(t-t_i)}_{\text{impactos do defeito}}
\;+\; \varepsilon(t),
\label{eq:dev-sinal}
\end{equation}

\noindent
com $f_r$ a rotação, $A_k$ as amplitudes harmônicas, $D$ a amplitude do defeito, $\zeta_d$
a taxa de amortecimento do impacto, $f_{\mathrm{res}}$ a ressonância excitada,
$u(\cdot)$ o degrau unitário e $\varepsilon(t)$ ruído gaussiano branco. Os instantes de
impacto são $t_i = i/\mathrm{BPFO} + \delta_i$, em que $\delta_i$ modela o escorregamento
das esferas. As frequências características vêm da geometria do rolamento, pelas expressões
da seção~\ref{sec:fund-rolamento}.

\subsection{Três cadeias comparadas sobre o mesmo sinal}

O mesmo sinal passou por três cadeias:

\begin{alineas}
  \item referência: amostragem a \SI{20}{\kilo\hertz}, que representa um acelerômetro de
        banda larga;
  \item com filtro: decimação para \SI{1}{\kilo\hertz} com filtro antirrebatimento;
  \item sem filtro: subamostragem para \SI{1}{\kilo\hertz} sem nenhum filtro.
\end{alineas}

No terceiro caso, uma componente acima da frequência de Nyquist é observada na frequência
aparente da Equação~\eqref{eq:fund-alias}.

Nenhuma das cadeias simuladas é idêntica à do dispositivo, que usa o filtro interno do
sensor com banda de \SI{184}{\hertz}. A segunda cadeia mantém conteúdo até
\SI{500}{\hertz}. A terceira é um caso extremo, usado para medir o rebatimento isolado,
porque o MPU6050 sempre opera com o filtro interno (Tabela~\ref{tab:fund-dlpf}).
\emandamento{Os experimentos ainda precisam ser refeitos com a banda de \SI{184}{\hertz}.}

Na comparação entre cadeias, os descritores foram a curtose, o fator de crista, o pico e o pico
a pico, calculados sobre \SI{2}{\second} de sinal. Os experimentos de limiar, métrica e partição
usam janelas de \SI{0,2}{\second}, com \SI{50}{\percent} de sobreposição, e outro conjunto: valor
eficaz, curtose, fator de crista, média absoluta, razão entre valor eficaz e média absoluta e fração
da energia por banda. A cadeia de \SI{20}{\kilo\hertz} tem quatro bandas e soma nove descritores, e
a de \SI{1}{\kilo\hertz} tem duas e soma sete.

O detector do \emph{firmware} usa quatro descritores: valor eficaz, curtose em excesso, fator de
crista e frequência dominante do espectro de velocidade. As fórmulas são iguais no treino e no
\emph{firmware}, com duas diferenças de sinal: o \emph{firmware} remove a média da janela e lê o eixo
Z, e o treino não remove a média e usa o canal radial do sensor \emph{underhang}. Nenhum desses
descritores mede a relação entre harmônicas, a fase ou a linha na frequência característica
(seção~\ref{sec:res-discussao}).

\subsection{Dado real decimado}

As taxas de detecção por modo de falha e a comparação entre limiares vieram das gravações do conjunto público
MAFAULDA~\cite{mafaulda}, obtidas em bancada instrumentada a \SI{50}{\kilo\hertz} e
decimadas em \emph{software} para \SI{1}{\kilo\hertz}. O conjunto soma \num{1951} arquivos
em seis classes (operação normal, desbalanceamento, desalinhamento horizontal e vertical e
defeito de rolamento nas posições \emph{overhang} e \emph{underhang}), segmentados em
\num{17559} janelas de \num{512} amostras, sem sobreposição. Apenas \num{49} arquivos são de operação normal,
e todo o treino saiu deles, porque o detector aprende somente a operação normal. A bancada
está abaixo do limite de potência e de altura de eixo da norma. Ela serve como fonte de
sinal para estudar a cadeia de medição, mas não como caso de aplicação da norma.

O limiar do detector foi calibrado pelo quantil do escore sobre arquivos normais não usados no
ajuste das árvores, com falso alarme alvo de \SI{1}{\percent}. O limiar padrão de \num{0,5} foi
avaliado nas mesmas partições.

A validação usou cinco partições por ensaio, com os arquivos de origem como grupos, para
que janelas de um mesmo ensaio não ficassem no treino e no teste ao mesmo tempo
(seção~\ref{sec:res-estimacao}). Nos ensaios de rolamento, os resultados foram separados
pelo nível de massa de desbalanceamento sobreposta, porque o valor agregado mistura os dois
efeitos.

\section{Conversão normativa e detecção}
\label{sec:dev-estimacao}

A norma define as zonas de severidade em velocidade eficaz de banda larga, e o acelerômetro
entrega aceleração \cite{iso208163_2022}. A conversão foi testada contra uma senoide cuja
integral é conhecida analiticamente, em cinco variantes: integração no tempo com e sem
filtro passa-alta, integração no domínio da frequência restrita à banda normativa, e esta
última e a integração no tempo sem filtro repetidas sob viés contínuo contido na
especificação do sensor. O método correto já está na literatura \cite{brandt2014}, e este
trabalho mediu o erro das alternativas em relação ao valor analítico. O ensaio usa
\SI{5}{\kilo\hertz} e \SI{4}{\second} de sinal. A integração no tempo é uma soma acumulada que parte
de zero, sem remoção da média, e o filtro passa-alta, quando usado, é de \SI{10}{\hertz}.

A detecção usou o Isolation Forest \cite{liu2008}, treinado apenas com dados de operação normal,
como no caso de uma máquina sem histórico rotulado. Em sinal sintetizado, na cadeia de
\SI{20}{\kilo\hertz}, o parâmetro de contaminação variou em cinco valores. O falso alarme foi medido
nas mesmas janelas do ajuste e, depois, em um ensaio normal deixado fora do treino, em cinco rodadas.

A métrica foi avaliada em defeito incipiente, com amplitude $D$ de \num{0,25}, contra \num{2,0} no
experimento de limiar. A área sob a curva ROC global e a área parcial \cite{koizumi2020} foram
comparadas entre as cadeias de \SI{20}{\kilo\hertz} e de \SI{1}{\kilo\hertz}, sobre as janelas
normais do ajuste e as janelas com defeito.

O protocolo de validação foi avaliado com um Random Forest supervisionado de \num{150} árvores,
treinado com as duas classes, na cadeia de \SI{1}{\kilo\hertz}. Esse classificador não é o detector
do \emph{firmware}.

\section{Dispositivo e \emph{firmware}}
\label{sec:dev-firmware}

O dispositivo é formado por um acelerômetro MEMS ligado por I2C a um microcontrolador ESP32
e por um módulo de rádio RA-02, com o circuito SX1278, ligado por SPI. A
Tabela~\ref{tab:dev-config} mostra a configuração usada nos ensaios. \emandamento{O sensor
da montagem de 24 de outubro de 2026 respondeu ao registro de identidade com o valor
\texttt{0x74}, do perfil MPU6515, e não com o \texttt{0x68} do MPU6050. A marcação do chip
ainda será conferida. As Tabelas~\ref{tab:fund-mpu6050} e~\ref{tab:fund-dlpf} são do MPU6050.}

\begin{table}[!htb]
\IBGEtab{%
\caption{Configuração do dispositivo nos ensaios de bancada}%
\label{tab:dev-config}%
}{%
\begin{tabular}{@{}ll@{}}
\toprule
Item & Valor \\
\midrule
Microcontrolador         & ESP32-S3 ou ESP32-C3 \\
Faixa do acelerômetro    & $\pm\SI{4}{g}$ \\
Filtro interno do sensor & \texttt{DLPF\_CFG}~=~1, banda de \SI{184}{\hertz} \\
Taxa de amostragem alvo  & \SI{1}{\kilo\hertz}, com tolerância de \SI{2}{\percent} \\
Janela                   & \num{512} amostras do eixo Z \\
Frequência do rádio      & \SI{433}{\mega\hertz} \\
Parâmetros do enlace     & SF9, BW de \SI{125}{\kilo\hertz}, CR de 4/5, preâmbulo de 8 símbolos \\
Potência de transmissão  & \SI{2}{dBm} \\
Pacote                   & \num{21} bytes, com CRC de 16 bits \\
Intervalo entre envios   & cerca de \SI{11}{\second} \\
\bottomrule
\end{tabular}%
}{%
\fonte{Elaboração própria.}%
}
\end{table}

O \emph{firmware} foi escrito em C++ com o \emph{framework} Arduino e compilado com o
PlatformIO. O mesmo código atende às duas placas, com a pinagem definida em um arquivo
próprio. Na bancada, cada
ciclo configura o acelerômetro, lê a janela pelo FIFO do sensor, calcula os quatro descritores,
envia o pacote pelo rádio e espera cerca de \SI{10}{\second} até o ciclo seguinte. O pacote leva
o identificador do dispositivo, o contador de ciclos, o estado, o código de erro, o valor eficaz e a
temperatura, que vai como NaN na bancada. Os outros três descritores só aparecem no registro serial.

Durante a leitura, o \emph{firmware} estima a taxa de amostragem real e descarta a janela
se ela se afastar mais de \SI{2}{\percent} de \SI{1}{\kilo\hertz}, porque a frequência
dominante é calculada com base nessa taxa.

Cada função devolve um código de estado, e o pacote leva o código do primeiro erro do ciclo.
Quando a leitura falha, a função não escreve no vetor de saída, o pacote sai com o estado
\texttt{UNKNOWN} e o valor eficaz vai como NaN. Uma versão anterior preenchia o vetor com
zeros em caso de erro, e os descritores calculados sobre esses zeros seriam classificados
como máquina sadia. O Apêndice~\ref{ape:praticas} descreve essas práticas.

O detector tem \num{100} árvores, treinadas em Python com subamostras de \num{256} pontos e
exportadas para um arquivo de cabeçalho em C++. Um teste compara o escore calculado em C++
com o da biblioteca usada no treino. \emandamento{A inferência ainda não é executada na
bancada: nesse perfil, o \emph{firmware} envia sempre o estado \texttt{UNKNOWN}.}

As funções de processamento não dependem do \emph{hardware} e rodam no computador. Os
\num{85} casos de teste da suíte passaram na versão avaliada, e parte deles cobre módulos fora da
montagem, como o termistor.

\section{Receptor}
\label{sec:dev-comunicacao}

O receptor é outro ESP32 com o mesmo módulo de rádio. Os parâmetros do enlace ficam em um
único arquivo, compartilhado com o transmissor, para que uma atualização da biblioteca de
rádio não os altere. Antes de aceitar um pacote, o receptor confere o tamanho, o CRC, a
versão e a validade dos campos. O dispositivo não tem relógio e envia apenas um contador de
ciclos, então é o receptor que informa o instante de chegada ao enviar os dados. A diferença entre
contadores consecutivos mostra quantos pacotes se perderam.

A implementação do CRC em C++ foi comparada com uma implementação em Python em \num{204} casos
de teste. \emandamento{O envio dos dados do receptor a um servidor por
Wi-Fi está implementado, mas não tem registro de ensaio.}

\section{Ensaio de bancada}
\label{sec:dev-bancada}

\emandamento{Foram analisados dois registros de bancada, salvos no repositório do projeto.
No primeiro, de 11 de outubro de 2026, um ESP32-S3 transmitiu e um ESP32-C3 recebeu, e o
registro é o do receptor. No segundo, de 24 de outubro de 2026, um ESP32-C3 transmitiu, e o
registro é o do transmissor, sem receptor associado. Os dois foram usados para avaliar os
critérios da hipótese H6: o desvio da taxa de amostragem em relação ao limite de
\SI{2}{\percent} e a validação dos pacotes pelo CRC. Esses registros serão substituídos
pelas medições com o dispositivo montado.}

%%%%%%%%%%%%%%%%%%%%%%%%%%%%%%%%%%%%%%%%%%%%%%%%%%%%%%%%%%%%%%%%%%%%%%%%%%%%%%%
%% Seção 4 --- Resultados e discussão
%%
%% Mesma ordem da seção de métodos: aquisição, estimação, ensaio de bancada e
%% comparação com a literatura. A origem de cada resultado (sinal sintetizado,
%% dado real decimado ou registro de bancada) é informada no próprio texto e nas
%% notas das tabelas e figuras.
%%
%% Fontes dos valores:
%%   sinal sintetizado    experiments/figures/data/fig1*.csv e fig2*.csv
%%                        (scripts/export_source_data.py)
%%   dado real decimado   experiments/figures/data/tab_deteccao_*.csv
%%                        (scripts/export_deteccao_data.py)
%%   bancada              logs/device-monitor-261011-155502.log e
%%                        logs/c3-device-validation-20261024-183820.log
%%%%%%%%%%%%%%%%%%%%%%%%%%%%%%%%%%%%%%%%%%%%%%%%%%%%%%%%%%%%%%%%%%%%%%%%%%%%%%%

\makeatletter
%% booktabs não é carregado pelo template; se ausente, as réguas viram \hline
\@ifundefined{toprule}{%
  \newcommand{\toprule}{\hline}%
  \newcommand{\midrule}{\hline}%
  \newcommand{\bottomrule}{\hline}%
}{}
\@ifundefined{emandamento}{%
  \@ifundefined{textcolor}%
    {\newcommand{\emandamento}[1]{#1}}%
    {\newcommand{\emandamento}[1]{\textcolor{red}{#1}}}%
}{}
\makeatother

%% Caminhos para as figuras geradas pelo pipeline do repositório.
%% NÃO copiar as figuras para dentro do TCC: elas são regeneráveis por
%% experiments/figures/ e uma segunda cópia divergiria em silêncio.
%% Ajustar o prefixo abaixo se a pasta de compilação mudar de lugar.
\providecommand{\kxfig}[1]{../../experiments/figures/output/#1}

\chapter{Resultados e discussão}
\label{cap:resultados}

Em dado real decimado, a detecção foi de \SI{75,0 \pm 3,5}{\percent} para desbalanceamento
e de no máximo \SI{20,7}{\percent} para defeito de rolamento sem massa de desbalanceamento
sobreposta. Com o limiar padrão, o falso alarme foi de \SI{31,2 \pm 4,2}{\percent}, contra
\SI{2,0 \pm 1,3}{\percent} com o limiar calibrado.

\section{Origem dos resultados}
\label{sec:proc}

Os resultados vêm de sinal sintetizado, de dado real decimado e de registros de bancada.

\textbf{Sinal sintetizado.} O sinal é gerado pelo modelo da seção~\ref{sec:dev-sinal}. Ele
        permite comparar métodos, mas não indica o desempenho em uma máquina.

\textbf{Dado real decimado.} As gravações vêm de uma bancada instrumentada e foram decimadas
        em \emph{software} para a taxa do dispositivo. O sinal é medido, mas a aquisição é
        simulada.

\textbf{Registro de bancada.} \emandamento{Os valores vêm do próprio dispositivo, sem motor
        e sem falha induzida. Eles permitem avaliar a aquisição e a transmissão, mas não a
        detecção.}

\section{Limites da aquisição}
\label{sec:res-aquisicao}

Com amostragem a \SI{1}{\kilo\hertz}, a ressonância em que aparece a energia do defeito de
rolamento fica fora da banda. Em dado real decimado, na condição sem massa sobreposta, o
desbalanceamento foi a única classe com detecção acima de \SI{20,7}{\percent}.

\subsection{Banda e evidência de defeito de rolamento}

A taxa de \SI{1}{\kilo\hertz} põe a frequência de Nyquist em \SI{500}{\hertz}. A frequência
característica de defeito de pista externa, de \SI{104,6}{\hertz} na geometria considerada,
fica dentro dessa banda, mas a ressonância de \SI{3,5}{\kilo\hertz} do modelo de sinal fica
fora.

Em sinal sintetizado, o contraste entre as condições com e sem defeito foi de \num{413}
vezes por análise de envelope na banda de ressonância, amostrada a \SI{20}{\kilo\hertz}, e
de \num{20} vezes por espectro direto a \SI{1}{\kilo\hertz} (Figura~\ref{fig:banda}). As
duas medidas diferem na banda e no método de demodulação, então esse par não permite dizer
quanto da perda vem da banda.

Sem filtro abaixo de $f_s/2$, a ressonância de \SI{3,5}{\kilo\hertz} reaparece rebatida em
\SI{500}{\hertz}, onde não há nenhuma componente física. A \SI{1}{\kilo\hertz} com filtro, a
curtose quase não mudou com o defeito, de \num{1,76} para \num{1,77}. Sem filtro, ela foi de
\num{1,82} para \num{4,55}, valor próximo do obtido a \SI{20}{\kilo\hertz}. Assim, um detector
pode separar as classes por causa do rebatimento e passar na validação. O MPU6050 sempre opera
com o filtro interno, mas o fabricante não informa a atenuação desse filtro acima da banda
(seção~\ref{sec:fund-mpu6050}), e este trabalho não a mediu.

\begin{figure}[!htb]
\caption{Efeito do limite de banda sobre a evidência de defeito de rolamento}
\label{fig:banda}
\centering
\includegraphics[width=\textwidth]{\kxfig{fig1_banda_sensor.pdf}}
\fonte{Elaboração própria.}
\legenda{Sinal sintetizado. \textbf{a}, Densidade espectral de potência para rolamento sadio e com defeito de pista externa, com eixo vertical em escala logarítmica. A faixa clara vai até \SI{500}{\hertz}, e a escura, até os \SI{184}{\hertz} do filtro usado no \emph{firmware}. \textbf{b}, Razão entre as amplitudes na frequência característica nas duas condições. \textbf{c}, Razão entre as condições para quatro descritores de impacto, calculados sobre \SI{2}{\second} de sinal; a linha tracejada marca a razão 1, sem separação. \textbf{d}, Frequência aparente da ressonância após subamostragem sem filtro.}
\end{figure}

\subsection{Detecção por modo de falha}

Em dado real decimado, com o limiar calibrado para falso alarme alvo de \SI{1}{\percent}, o
falso alarme medido foi de \SI{2,0 \pm 1,3}{\percent} (média e desvio-padrão entre as cinco
partições). A Tabela~\ref{tab:deteccao} mostra a detecção por classe nesse ponto de
operação.

\begin{table}[!htb]
\IBGEtab{%
\caption{Detecção por modo de falha a \SI{2}{\percent} de falso alarme medido, em dado real decimado}%
\label{tab:deteccao}%
}{%
\begin{tabular}{lc}
\toprule
Modo de falha & Detecção (\si{\percent}) \\
\midrule
Desbalanceamento                      & \num{75,0 \pm 3,5} \\
Rolamento \emph{overhang} (puro)      & \num{20,7 \pm 7,0} \\
Rolamento \emph{underhang} (puro)     & \num{14,8 \pm 4,3} \\
Desalinhamento vertical               & \num{7,2 \pm 4,3} \\
Desalinhamento horizontal             & \num{5,0 \pm 3,1} \\
\bottomrule
\end{tabular}%
}{%
\fonte{Elaboração própria, com dados de \citeonline{mafaulda}.}%
\nota{A incerteza é o desvio-padrão entre as cinco partições. As duas classes de rolamento correspondem à condição sem massa de desbalanceamento sobreposta.}%
}
\end{table}

Nas duas classes de desalinhamento, o desvio entre partições é da ordem do próprio valor, e
\SI{7,2 \pm 4,3}{\percent} admite valores próximos de zero. Por isso, conclusões baseadas
nesses dois números são fracas (seção~\ref{sec:res-discussao}).

As classes de rolamento foram ensaiadas com \SI{0}{\gram}, \SI{6}{\gram}, \SI{20}{\gram} e \SI{35}{\gram} de massa de
desbalanceamento sobreposta, e a detecção foi maior nas massas maiores. Na classe
\emph{underhang}, ela passou de \SI{14,8}{\percent} sem massa para \SI{87,3}{\percent} com
\SI{35}{\gram}. Na classe \emph{overhang}, passou de \SI{20,7}{\percent} para
\SI{78,3}{\percent}. Os agregados das quatro condições são de \SI{55,9}{\percent} e
\SI{42,2}{\percent}, mas esses números não podem ser lidos como detecção de defeito de
rolamento, porque o modelo parece responder à massa adicionada.

\section{Decisões de estimação}
\label{sec:res-estimacao}

Quatro decisões alteraram os resultados: o método de integração, o limiar, a métrica e o
protocolo de validação (Figura~\ref{fig:metodologia}). Exceto onde indicado, os valores desta
seção vêm de sinal sintetizado.

\begin{figure}[!htb]
\caption{Efeito de quatro escolhas metodológicas sobre os resultados}
\label{fig:metodologia}
\centering
\includegraphics[width=\textwidth]{\kxfig{fig2_metodologia.pdf}}
\fonte{Elaboração própria.}
\legenda{Sinal sintetizado. \textbf{a}, Velocidade eficaz por cinco variantes de integração, comparada com o valor analítico (linha tracejada), com soma acumulada sem remoção da média nas variantes no tempo. \textbf{b}, Velocidade sob viés de \SI{0,02}{\meter\per\second\squared}. \textbf{c}, Acurácia de um Random Forest supervisionado sobre os mesmos dados, variando a forma de partição; as barras de erro são o desvio-padrão ($n = 10$ sorteios e $n = 5$ partições), e os pontos, as partições individuais. \textbf{d}, Fração das janelas normais do ajuste classificada como anomalia, por valor do parâmetro de contaminação. \textbf{e}, Curvas ROC (característica de operação do receptor) em defeito incipiente; a faixa marca falso alarme de até \SI{10}{\percent}, usado na área parcial (pAUC).}
\end{figure}

\subsection{Integração de aceleração para velocidade}

A norma define as zonas de severidade em velocidade eficaz, e o acelerômetro entrega
aceleração, então o sinal precisa ser integrado. A Tabela~\ref{tab:res-integracao} compara os
métodos contra uma senoide de \SI{3,0}{\meter\per\second\squared} a \SI{50}{\hertz}, cuja
integral é conhecida e tem valor eficaz de \SI{6,7524}{\milli\meter\per\second}.

\begin{table}[!htb]
\IBGEtab{%
\caption{Velocidade eficaz por variante de integração, em relação ao valor analítico}%
\label{tab:res-integracao}%
}{%
\begin{tabular}{lcc}
\toprule
Método & Resultado (\si{\milli\meter\per\second}) & Erro (\si{\percent}) \\
\midrule
Soma acumulada, sem filtro passa-alta               & \num{11,69} & $+\num{73}$ \\
Soma acumulada, com filtro passa-alta               & \num{17,31} & $+\num{156}$ \\
\textbf{Integração na frequência, banda normativa}  & \num{6,75}  & \textbf{\num{0}} \\
Na frequência, com viés de \SI{0,02}{\meter\per\second\squared} & \num{6,75} & \num{0} \\
Soma acumulada, sem filtro, com o mesmo viés        & \num{55,08} & $+\num{716}$ \\
\bottomrule
\end{tabular}%
}{%
\fonte{Elaboração própria.}%
\nota{Sinal sintetizado, a \SI{5}{\kilo\hertz} durante \SI{4}{\second}. Erro em relação ao valor analítico, arredondado para inteiro. A soma acumulada parte de zero, sem remoção da média.}%
}
\end{table}

A soma acumulada parte de zero e deixa a constante de integração na média da velocidade, de
\SI{9,55}{\milli\meter\per\second}. Daí vem o erro: $\sqrt{\num{6,75}^2 + \num{9,55}^2} =
\SI{11,69}{\milli\meter\per\second}$. Com a média removida, o erro cai para
\SI{0,02}{\percent} sem filtro e para \SI{1,6}{\percent} com filtro. Um viés de
\SI{0,02}{\meter\per\second\squared}, contido na especificação do sensor, levou a soma a
\SI{55,08}{\milli\meter\per\second}, e a remoção da tendência linear também trouxe esse valor
de volta a \SI{6,75}{\milli\meter\per\second}. Sem esse cuidado, o sinal seria classificado na
zona D em qualquer máquina, enquanto o valor correto fica abaixo do limite entre as zonas C e
D em três das quatro combinações da norma.

A integração na frequência deu \SI{6,7525}{\milli\meter\per\second}, a menos de
\SI{0,01}{\percent} do valor analítico, com e sem o viés, porque o termo contínuo fica fora da
banda normativa. Nesse ensaio, os \SI{50}{\hertz} coincidem com um coeficiente exato da DFT, e
os \SI{4}{\second} de registro são mais longos que a janela de \SI{0,512}{\second} do
dispositivo.

\subsection{Limiar e falso alarme}

Em dado real decimado, avaliado em arquivos fora do treino, o limiar padrão de \num{0,5} marcou
\SI{31,2 \pm 4,2}{\percent} da operação normal como anomalia, e o limiar calibrado marcou
\SI{2,0 \pm 1,3}{\percent}.

Em sinal sintetizado, na cadeia de \SI{20}{\kilo\hertz}, o valor padrão do parâmetro de
contaminação marcou \SI{42,1}{\percent} das \num{140} janelas normais, e o valor \num{0,01}
marcou \SI{1,4}{\percent}. Essas taxas foram medidas nas próprias janelas do ajuste. Quando um
ensaio normal ficou fora do treino, os dois valores marcaram \SI{100}{\percent} das janelas nas
cinco rodadas. Os ensaios sintetizados têm ganhos diferentes, e o limiar ajustado em quatro deles
não serviu para o quinto.

Nos cinco valores testados, os escores foram idênticos, e só o deslocamento do limiar mudou. A
detecção ficou em \SI{100}{\percent} e a área sob a curva ROC em \num{1,000}, porque as classes
estão perfeitamente separadas nesse sinal, o que é improvável em dado real. O valor de
\SI{42}{\percent} vale só para este sinal.

\subsection{Métrica global e métrica no ponto de operação}

Em defeito incipiente, a área sob a curva completa caiu de \num{1,000} para \num{0,963} entre as
cadeias de \SI{20}{\kilo\hertz} e de \SI{1}{\kilo\hertz}, uma queda de \SI{3,7}{\percent}. A área
parcial caiu de \num{0,999} para \num{0,837} com limite de \SI{10}{\percent} de falso alarme, uma
queda de \SI{16,2}{\percent}, e de \num{0,998} para \num{0,795} com limite de \SI{5}{\percent}, uma
queda de \SI{20,3}{\percent}. Na região de baixo falso alarme, a queda relativa foi de \num{4,4} a
\num{5,5} vezes a indicada pela área global.

A cadeia de \SI{20}{\kilo\hertz} usa nove descritores e a de
\SI{1}{\kilo\hertz} usa sete, então a queda mistura o efeito da banda e o dos descritores. As áreas
foram calculadas sobre as janelas normais do ajuste e as janelas com defeito. Pela forma do desafio
de referência, dividida só por $\beta$, as quedas seriam de \SI{30,8}{\percent} e
\SI{39,7}{\percent}, o que muda os valores, mas não o sentido da comparação.

\subsection{Partição por janela e partição por ensaio}

Um Random Forest supervisionado, sobre os mesmos dados e os mesmos descritores, atingiu
\SI{99,0 \pm 1,0}{\percent} de acurácia sob partição aleatória por janela e
\SI{63,2 \pm 19,2}{\percent} sob partição por ensaio, uma diferença de aproximadamente \num{36}
pontos percentuais. Na partição por ensaio, três das cinco partições ficaram próximas de
\SI{50}{\percent}, o nível do acerto ao acaso, e uma alcançou \SI{100}{\percent}. As janelas têm
\SI{50}{\percent} de sobreposição, o que favorece a partição aleatória, porque janelas quase iguais
ficam no treino e no teste.

\section{Ensaio de bancada}
\label{sec:res-bancada}

\emandamento{No registro de 11 de outubro de 2026, com um ESP32-S3 como transmissor e um ESP32-C3
como receptor, o receptor aceitou \num{270} pacotes em \SI{62}{\minute}, todos com CRC de 16 bits
correto. Houve um erro de CRC no próprio rádio, que corresponde ao único pacote perdido. A contagem de
sequência recomeçou em 1 uma vez, e o receptor ficou \SI{13,4}{\minute} sem pacotes nesse intervalo. A
intensidade do sinal recebido ficou entre \SI{-73}{dBm} e \SI{-55}{dBm}, e a relação sinal-ruído,
entre \SI{9,5}{\decibel} e \SI{14}{\decibel}. Antes do primeiro pacote, o receptor registrou uma
instabilidade no barramento SPI e reiniciou o rádio. A distância entre os módulos não foi registrada.}

\emandamento{No registro de 24 de outubro de 2026, com um ESP32-C3 como transmissor, os três ciclos
gravados terminaram sem erro. A taxa de amostragem estimada ficou entre \SI{999,78}{\hertz} e
\SI{999,88}{\hertz}, com desvio máximo de \SI{0,022}{\percent}, abaixo do limite de
\SI{2}{\percent}. O intervalo entre transmissões foi de \SI{11,0}{\second}, e a fase de transmissão
durou cerca de \SI{195}{\milli\second}, próximo dos \SI{185}{\milli\second} calculados na
seção~\ref{sec:fund-lora}.}

\emandamento{Os dois registros vêm de montagens diferentes, e o segundo tem só três ciclos. Eles
indicam que a aquisição e a transmissão funcionam em bancada, mas ainda não permitem avaliar H6 por
completo. A escolha final do fator de espalhamento depende da distância até o receptor, que ainda não
foi medida.}

\section{Comparação com a literatura}
\label{sec:res-discussao}

A diferença de \num{36} pontos percentuais entre as duas formas de partição tem correlato na
literatura, que mede diferenças da ordem de dezenas de pontos percentuais em uma base pública de
rolamentos~\cite{varejao2025,vieira2026}. A concordância é de direção e de ordem de grandeza, porque
o valor deste trabalho vem de sinal sintetizado e de um classificador supervisionado.

Outros autores também separam desbalanceamento e rolamento no projeto do nó. Um nó publicado usa um
acelerômetro de banda estreita para desbalanceamento e outro de dezenas de quilohertz para
rolamento~\cite{jakobsen2024}, e o trabalho mais próximo em arquitetura demonstra desbalanceamento e
deixa rolamento para depois~\cite{kolok2025}.

O falso alarme de \SI{31,2}{\percent} com o limiar padrão, em dado real, fica próximo das
implementações embarcadas publicadas, que relatam de \SI{22}{\percent} a \SI{38}{\percent} sobre
operação normal~\cite{chevtchenko2023,kolok2025}. Esses trabalhos usam conjuntos, máquinas e limiares
diferentes, e a comparação serve só como ordem de grandeza.

Os resultados de detecção não mostram que a banda seja a causa da baixa detecção de rolamento. O
detector marca afastamento da operação normal e não identifica o tipo de defeito. Os quatro descritores
não medem a linha na frequência característica, embora o espectro direto a \SI{1}{\kilo\hertz} tenha
mostrado contraste de \num{20} vezes nessa linha em sinal sintetizado. O desalinhamento, que aparece
dentro da banda, teve detecção menor (\SI{5,0}{\percent} e \SI{7,2}{\percent}) que o defeito de
rolamento (\SI{14,8}{\percent} e \SI{20,7}{\percent}). Além disso, todas as classes usam o mesmo canal
da bancada, o sensor radial do lado \emph{underhang}, mesmo no defeito \emph{overhang}.

A baixa detecção de desalinhamento pode vir da escolha dos descritores. Sua assinatura é a razão entre
as linhas de duas vezes e uma vez a rotação, que valor eficaz, curtose, fator de crista e frequência
dominante não medem. A evidência é fraca, porque os valores têm dispersão da ordem do próprio valor. A
seção~\ref{sec:conc-futuro} indica o experimento que testaria essa explicação.

%%%%%%%%%%%%%%%%%%%%%%%%%%%%%%%%%%%%%%%%%%%%%%%%%%%%%%%%%%%%%%%%%%%%%%%%%%%%%%%
%% Seção 5 --- Conclusões
%%
%% Volta ao problema da Introdução, avalia as hipóteses H1 a H6 e não traz autores
%% nem citações novas, como pede o modelo do iCEV. Ao alterar a resposta ao problema,
%% alterar também o resumo e o abstract.
%%
%% MARCAÇÃO DE TRECHO EM ANDAMENTO --------------------------------------------
%% \emandamento{...} destaca em vermelho os trechos que mudam com as medições do
%% dispositivo montado. Ver nota no capitulo01.
%%%%%%%%%%%%%%%%%%%%%%%%%%%%%%%%%%%%%%%%%%%%%%%%%%%%%%%%%%%%%%%%%%%%%%%%%%%%%%%

\makeatletter
\@ifundefined{emandamento}{%
  \@ifundefined{textcolor}%
    {\newcommand{\emandamento}[1]{#1}}%
    {\newcommand{\emandamento}[1]{\textcolor{red}{#1}}}%
}{}
\makeatother

\chapter{Conclusões}
\label{cap:conclusoes}

Este trabalho perguntou quais modos de falha de motores uma cadeia de medição de baixo custo
permite detectar e quanto as decisões de estimação alteram o desempenho reportado. Em dado
real decimado para a taxa do dispositivo, com falso alarme medido de \SI{2,0 \pm 1,3}{\percent},
a detecção foi de \SI{75,0 \pm 3,5}{\percent} em desbalanceamento, de \SI{14,8 \pm 4,3}{\percent}
e \SI{20,7 \pm 7,0}{\percent} em defeito de rolamento puro, e de \SI{5,0 \pm 3,1}{\percent} e
\SI{7,2 \pm 4,3}{\percent} em desalinhamento horizontal e vertical. Com esse detector e esses
quatro descritores, a cadeia pode ser usada para detectar desbalanceamento, mas detecta pouco
defeito de rolamento sem desbalanceamento sobreposto.

As decisões de estimação mudaram os resultados. O limiar padrão marcou \SI{31,2 \pm
4,2}{\percent} da operação normal em dado real, contra \SI{2,0 \pm 1,3}{\percent} do limiar
calibrado. Para comparar dispositivos dessa classe, é preciso informar a banda, o limiar e o
protocolo de validação.

\section{Avaliação das hipóteses}
\label{sec:conc-hipoteses}

As hipóteses H1, H2, H4 e H5 não foram refutadas. A H3 não foi refutada em dado real, mas foi
refutada em um teste com sinal sintetizado. A H6 depende do dispositivo e ainda está em avaliação.

\textbf{H1, banda.} Não foi refutada. A detecção de desbalanceamento superou a de defeito de
        rolamento sem massa sobreposta, sob a mesma taxa de falso alarme. Esse resultado não
        isola o efeito da banda, porque o desalinhamento, dentro da banda, teve detecção ainda
        menor, e os descritores não medem a linha na frequência característica.
        \emandamento{Os experimentos mantêm conteúdo até \SI{500}{\hertz}, e o \emph{firmware} usa
        \SI{184}{\hertz}. A avaliação precisa ser refeita com essa banda.}

\textbf{H2, integração.} Não foi refutada. Contra uma senoide com integral conhecida, a soma
        acumulada errou $+\SI{73}{\percent}$ sem filtro passa-alta, $+\SI{156}{\percent}$ com ele
        e $+\SI{716}{\percent}$ sob viés, e a integração na frequência ficou a menos de
        \SI{0,01}{\percent} do valor analítico. Quase todo o erro da soma vem da constante de
        integração: com a média removida, ele cai para \SI{0,02}{\percent} sem filtro e para
        \SI{1,6}{\percent} com filtro.

\textbf{H3, limiar.} Em dado real, fora do treino, não foi refutada: o limiar padrão marcou
        \SI{31,2 \pm 4,2}{\percent} da operação normal, e o calibrado, \SI{2,0 \pm 1,3}{\percent}.
        Em sinal sintetizado, nas janelas do ajuste, os valores foram de \SI{42,1}{\percent} e
        \SI{1,4}{\percent}. Com um ensaio normal fora do treino, os dois marcaram
        \SI{100}{\percent}, o que refuta a hipótese nessa condição. Os escores foram idênticos nos
        cinco valores do parâmetro.

\textbf{H4, métrica.} Não foi refutada. Entre as cadeias de \SI{20}{\kilo\hertz} e de
        \SI{1}{\kilo\hertz}, a área global caiu \SI{3,7}{\percent}, e a área parcial caiu
        \SI{16,2}{\percent} com limite de \SI{10}{\percent} de falso alarme e \SI{20,3}{\percent}
        com limite de \SI{5}{\percent}. As duas cadeias usam conjuntos diferentes de descritores,
        e a queda não pode ser atribuída só à banda.

\textbf{H5, protocolo.} Não foi refutada. Um Random Forest supervisionado atingiu
        \SI{99,0 \pm 1,0}{\percent} de acurácia sob partição aleatória por janela e
        \SI{63,2 \pm 19,2}{\percent} sob partição por ensaio. O experimento não usou o detector do
        \emph{firmware}.

\textbf{H6, dispositivo.} \emandamento{Ainda em avaliação. Em bancada, o receptor aceitou \num{270}
        pacotes com CRC correto, e um pacote se perdeu. Em outra montagem, a taxa de amostragem ficou
        entre \SI{999,78}{\hertz} e \SI{999,88}{\hertz} em três ciclos, dentro do limite de
        \SI{2}{\percent}. Os critérios precisam ser repetidos na montagem final, com transmissor e
        receptor registrados no mesmo ensaio.}

\section{Limites dos resultados}
\label{sec:conc-alcance}

Os resultados de detecção vêm de dado real de uma bancada instrumentada, decimado em \emph{software}.
O sinal foi medido, mas a aquisição do dispositivo foi simulada. O treino usou apenas \num{49} arquivos
de operação normal, e todas as classes usam o mesmo canal da bancada. Essa bancada está abaixo do
limite de potência e de altura de eixo da norma, então os valores não podem ser estendidos diretamente
a máquinas industriais.

Os resultados de banda, integração, métrica e protocolo vêm de sinais sintetizados. Seus valores
absolutos dependem da amplitude do defeito, da frequência de ressonância e do nível de ruído escolhidos,
e devem ser lidos como ordens de grandeza. As relações entre os métodos devem depender menos desses
parâmetros, mas este trabalho não fez a varredura que confirmaria isso.

\emandamento{O dispositivo só foi ensaiado em bancada, sem motor e sem falha induzida, e a inferência
ainda não é executada nesse perfil. Portanto, os valores de detecção não são do dispositivo em operação.}

\section{Trabalho futuro}
\label{sec:conc-futuro}

\textbf{Projeto de \emph{hardware} próprio.} Este trabalho usou componentes em montagem de bancada. As
placas de circuito impresso, o invólucro, a bateria e o comportamento térmico ainda precisam ser
projetados e ensaiados. O invólucro e a fixação acrescentam uma ressonância entre o motor e o sensor, e
essa frequência deve ser medida por excitação de impacto.

\textbf{Ensaio em motor.} \emandamento{O dispositivo precisa ser ensaiado em um motor, com a inferência
ativada e com falhas induzidas. O piso de ruído também deve ser medido, para separar um defeito não
detectado de um defeito abaixo do ruído.}

\textbf{Varredura de taxa de amostragem.} A mesma gravação, decimada em várias taxas e analisada por um
único método e um único conjunto de descritores, permitiria separar o efeito da banda dos demais.

\textbf{Efeito do filtro em dado medido.} Em sinal sintetizado, o rebatimento mudou a curtose. Falta
medir o mesmo efeito em uma gravação real, com e sem filtro.

\textbf{Limiar em dado industrial.} A varredura do parâmetro de contaminação contra o falso alarme medido,
em vibração industrial de banda limitada, permitiria verificar se a calibração por risco alvo supera o
valor padrão.

\textbf{Treino no próprio microcontrolador.} A busca não encontrou taxa de falso alarme publicada para esse
caso, que é o de uma máquina sem histórico rotulado.

\textbf{Descritores de razão harmônica.} A detecção de desalinhamento foi a menor entre os modos medidos, e
os descritores adotados não medem a relação entre harmônicas. O teste pode usar as mesmas gravações
decimadas.


%%%%%%%%%%%%%%%%%%%%%%%%%%%%%%%%%%%%%%%%%%%%%%%%%%%%%%%%%%%%%%%%%%%%%%%%%%%%%%%
%% ELEMENTOS PÓS-TEXTUAIS
%%%%%%%%%%%%%%%%%%%%%%%%%%%%%%%%%%%%%%%%%%%%%%%%%%%%%%%%%%%%%%%%%%%%%%%%%%%%%%%
\postextual

\citeoption{abnt-emphasize=bf}
\bibliography{referencias}

\begin{apendicesenv}
  %% "APÊNDICE A --- TÍTULO" (NBR 14724, 4.2.3.3), com a mesma separação das demais seções.
  %% Sem espaço antes do traço: a classe já deixa um depois da letra.
  \renewcommand{\ABNTEXcaptiondelim}{\textemdash~}
  \setlength{\afterchapskip}{\baselineskip}
  %%%%%%%%%%%%%%%%%%%%%%%%%%%%%%%%%%%%%%%%%%%%%%%%%%%%%%%%%%%%%%%%%%%%%%%%%%%%%%%
%% Apêndice A --- Práticas de engenharia do firmware
%%
%% Valores conferidos em src/sensors/vibration.h, src/comms/lora.h,
%% src/ml/model.h, src/power/watchdog.h e training/tests/test_export_cpp.py.
%%%%%%%%%%%%%%%%%%%%%%%%%%%%%%%%%%%%%%%%%%%%%%%%%%%%%%%%%%%%%%%%%%%%%%%%%%%%%%%

\chapter{Práticas de engenharia do \emph{firmware}}
\label{ape:praticas}

Algumas práticas usadas no \emph{firmware} vêm de normas e da literatura de sistemas críticos.
Essas fontes servem aqui como referência. O Kaelix é um protótipo acadêmico e não passou por
nenhum processo de certificação.

\section{Fontes das práticas}
\label{sec:ape-escalas}

As normas setoriais classificam a criticidade de formas diferentes, e não há equivalência
normativa entre elas. A norma de aviônica classifica o \emph{software} pela severidade da
condição de falha e não atribui taxa de falha ao \emph{software} \cite{do178c}. A norma
genérica define faixas numéricas de medida de falha alvo \cite{iec61508}. A norma automotiva
combina severidade, exposição e controlabilidade \cite{iso26262}.

\section{Práticas adotadas}
\label{sec:ape-mecanismos}

Nenhuma das práticas abaixo exige componente adicional.

\textbf{Propagação de erro.} Cada função devolve um código de estado, e o pacote leva o código do
primeiro erro do ciclo. Sem esse código, uma leitura com falha chegaria ao receptor como uma leitura
normal \cite{esa1996,knight1986}.

\textbf{Saída só em caso de sucesso.} Quando a leitura do sensor falha, a função não escreve no
vetor de amostras. Assim, o detector não recebe zeros que pareçam uma máquina parada.

\textbf{Valor inválido explícito.} O valor eficaz vai como NaN quando não foi medido, porque zero
é uma leitura possível. O estado \texttt{UNKNOWN} indica que o dispositivo não avaliou a máquina.

\textbf{Versão no pacote.} O primeiro byte do pacote informa a versão do formato, e o receptor
recusa versões que não conhece.

\textbf{\emph{Watchdog}.} Um temporizador reinicia o dispositivo se um ciclo travar.

\textbf{Testes no computador.} As funções de processamento rodam sem o \emph{hardware}. Um teste
compara o escore do detector em C++ com o da biblioteca usada no treino, e outro compara o CRC
em C++ com uma implementação em Python.

\section{Limites das práticas}
\label{sec:ape-limites}

Essas práticas não substituem redundância de \emph{hardware}, e a suposição de independência
entre versões diversas de \emph{software} não se confirmou empiricamente \cite{knight1986}. Os
testes no computador não cobrem o comportamento do rádio nem o do sensor.
\emandamento{Nenhum registro de bancada provoca uma falha de propósito para testar a propagação
de erro no dispositivo.}

\end{apendicesenv}

\end{document}
